\chapter{基于数据本身的知识获取}
\label{chap:data-derived-knowledge-acquisition}

我们已经有一批花卉照片，却还没有逐张确认物种。照片中的花瓣、叶片和拍摄背景仍然可以被读取，同一张照片也可以裁成两个视图；缺少类别答案，并不等于没有任何可用于学习的信息。前几章取得的类别、数值或成对判断直接回答既定任务，本章则从这些尚未附带任务答案的记录出发：怎样把已有内容和结构变成可训练、可检查的依据？

先在同一批照片上比较三种做法。遮住一块花瓣后预测其像素，答案可以在遮蔽前保留；让一个裁剪预测另一个裁剪的表示，答案则要等编码器把图像转成特征向量以后才出现；先把照片聚成群组，再预测组号，答案又取决于分组算法。三种做法都使用预测，却没有取得同一种知识。沿用第~\ref{chap:knowledge-acquisition-overview}章的区分，构造出的目标或关系先是候选知识，进入损失计算后才承担监督信号的作用；能够构造它们，还不能确认它们对物种识别有用。

由目标的取得依据，我们可以分出三条路线。\term{观测内容}（observed content）是已经记录的词元、像素或状态，从中保留一部分便能得到内容答案；\term{样本关系}（sample relations）规定哪些记录应保留共同信息、哪些记录允许相对区分，再据此约束表示；\term{数据分组}（data grouping）把多个样本组织成群组，用组号或对代表向量的分配作为临时答案。上下文和时间顺序限制内容怎样切分，多模态配对规定不同记录怎样对应，它们没有另添一种目标来源。沿这三条路线，我们既要追踪已有信息怎样进入目标和参数更新，也要检查训练后保留什么表示供下游使用。

图~\ref{fig:data-derived-knowledge-target-sources}让我们先沿两条路径找答案。左侧目标直接从原记录分出，模型只需计算预测；右侧两路都经过编码分支，比较的是训练当步算出的向量，参数改变后向量也会改变。图只回答目标来自哪里，还没有规定两路怎样更新。第~\ref{sec:data-derived-knowledge-training-objectives}节会据此区分两路共同求导与把一路输出暂时当作常量的安排。

\begin{figure*}[htbp]
  \centering
  \begin{tikzpicture}[
  x=1mm,y=1mm,
  font=\figurefont\footnotesize,
  panel/.style={draw=BookRule,fill=NNPanel,line width=.85pt,rounded corners=1.5mm},
  data/.style={draw=FigureInput,fill=FigureInput!12,line width=1.1pt,minimum height=10mm,align=center,inner sep=2mm},
  model/.style={draw=FigureModel,fill=FigureModel!12,line width=1.1pt,minimum height=10mm,align=center,inner sep=2mm},
  target/.style={draw=FigureOutput,fill=FigureOutput!10,line width=1.1pt,minimum height=10mm,align=center,inner sep=2mm},
  loss/.style={draw=FigureLoss,fill=FigureLoss!12,line width=1.1pt,minimum height=10mm,align=center,inner sep=2mm},
  flow/.style={knowledge arrow,draw=NNWire,line width=1.2pt,rounded corners=1mm},
  keep/.style={knowledge arrow,draw=FigureOutput,line width=1.1pt,rounded corners=1mm},
  annotation/.style={font=\figurefont\scriptsize,text=BookMuted,align=center}
]
  \draw[panel] (1,-76) rectangle (81,12);
  \draw[panel] (87,-76) rectangle (168,12);
  \node[font=\figurefont\bfseries\footnotesize,text=BookInk] at (41,7) {(a) 观测目标：答案随观测保留};
  \node[data,minimum width=24mm] (obs) at (41,-9) {完整观测\\$\bm{x}$};
  \node[data,minimum width=24mm] (visible) at (20,-28) {可见预测输入\\$\bm{x}_{V}$};
  \node[target,minimum width=24mm] (held) at (62,-28) {保留目标\\$\bm{x}_{M}$};
  \node[model,minimum width=24mm] (predict) at (20,-48) {模型预测\\$\widehat{\bm{x}}_{M}$};
  \node[loss,minimum width=24mm] (recloss) at (62,-48) {重建损失};
  \coordinate (obs-split) at (41,-18);
  \draw[NNWire,line width=1.2pt] (obs.south) -- (obs-split);
  \draw[flow] (obs-split) -| (visible.north);
  \draw[keep] (obs-split) -| (held.north);
  \draw[flow] (visible.south) -- (predict.north);
  \draw[flow] (predict.east) -- (recloss.west);
  \draw[keep] (held.south) -- (recloss.north);
  \node[annotation,text width=68mm] at (41,-66) {目标从原始观测中预先保留，\\不需要编码分支生成答案。};

  \node[font=\figurefont\bfseries\footnotesize,text=BookInk] at (127.5,7) {(b) 表示目标：答案在训练中计算};
  \node[data,minimum width=24mm] (raw) at (127.5,-9) {原样本\\$\bm{x}$};
  \node[data,minimum width=24mm] (viewa) at (106,-28) {视图A\\$\bm{v}$};
  \node[data,minimum width=24mm] (viewb) at (149,-28) {视图B\\$\bm{v}'$};
  \node[model,minimum width=24mm] (enca) at (106,-48) {编码分支\\$\bm{z}$};
  \node[model,minimum width=24mm] (encb) at (149,-48) {编码分支\\$\bm{z}'$};
  \node[loss,minimum width=22mm] (match) at (127.5,-66) {匹配目标};
  \coordinate (view-split) at (127.5,-18);
  \draw[NNWire,line width=1.2pt] (raw.south) -- (view-split);
  \draw[flow] (view-split) -| (viewa.north);
  \draw[flow] (view-split) -| (viewb.north);
  \draw[flow] (viewa.south) -- (enca.north);
  \draw[flow] (viewb.south) -- (encb.north);
  \draw[flow] (enca.south) |- (match.west);
  \draw[flow] (encb.south) |- (match.east);

  \node[annotation,anchor=north west,text width=166mm,align=left] at (1,-80) {右侧表示在编码之后产生，随模型参数和训练步变化；本图只比较目标来源，不预设教师或停止梯度。};
\end{tikzpicture}

  \caption{两种目标在何处产生。左侧从完整记录分出答案，只让可见部分进入预测模型；右侧从对应视图分别计算表示，再交给匹配目标。右侧画出的是前向计算，不表示已经指定独立教师或停止梯度。}
  \label{fig:data-derived-knowledge-target-sources}
\end{figure*}
\FloatBarrier

\section{基于观测内容的知识获取}
\label{sec:data-derived-knowledge-observation}

对一张尚未标注物种的照片，我们可以先保留花瓣区域的像素，再从模型输入中遮住该区域。模型不知道答案，训练过程却有原记录可供核对。\term{观测目标}（observation target）指这样从已记录内容中取出的预测对象，词元、属性和未来片段也可以承担同一角色。目标不需要正在学习的编码器生成，但哪些内容成为答案、哪些内容允许被读取，仍由遮蔽位置、预测方向和预处理规则决定。

为了检查答案有没有混进输入，需要把一条记录的几个角色分开。记完整观测为$\bm{x}=(x_j)_{j\in I}$，$I$枚举词元、像素块或字段等位置，$x_j$表示该位置的内容。构造规则从$I$中选择可见集合$V\subset I$与目标集合$M\subseteq I$，通常限制输入不能直接给出$\bm{x}_M$的原值。模型读取处理后的$\widetilde{\bm{x}}_V$，预测$\widehat{\bm{x}}_M$或相应条件分布。保存完整观测、处理后输入、保留目标及损失位置，才能逐步核对预处理是否把答案带回输入；把它们合存为一份未区分角色的“训练样本”，会失去这条检查路径。

表~\ref{tab:data-derived-knowledge-observation-tasks}按输入怎样改变、答案从哪里保留、误差在哪些位置计算来比较四种安排。同一未来片段也可以被扰动后用于恢复，所以四行不是互斥任务。决定这些安排差别的是可见信息和目标范围，而不是最后是否使用同一个损失名称。

\begin{table*}[htbp]
  \centering
  \renewcommand{\arraystretch}{1.15}
  \begin{tabularx}{\linewidth}{@{}>{\RaggedRight\setlength{\parindent}{0pt}\arraybackslash}p{22mm}*{3}{>{\RaggedRight\setlength{\parindent}{0pt}\arraybackslash}X}@{}}
    \toprule
    构造方式 & 模型可见部分 & 保留目标 & 损失位置 \\
    \midrule
    掩码恢复 & 隐藏选中图像块后的图像 & 被遮蔽块的原像素 & 通常只在掩码位置 \\
    去噪恢复 & 加入扰动后的图像 & 扰动前的原始记录 & 受扰位置或全部位置 \\
    上下文预测 & 前缀“花朵逐渐” & 后续词元“开放” & 约定的预测位置 \\
    时序预测 & 判断时点以前的历史窗口 & 后续帧或状态 & 一个未来点或一段轨迹 \\
    \bottomrule
  \end{tabularx}
  \caption{按可见信息、保留答案和计分位置比较观测目标。未来片段虽已进入训练记录，在真正预测的时点却不可读取；四种安排可组合，不能据表格分行认定它们互斥。}
  \label{tab:data-derived-knowledge-observation-tasks}
\end{table*}
\FloatBarrier

\subsection{基于遮蔽恢复的目标构造}
\label{sec:data-derived-knowledge-masking-denoising}

先固定一张完整照片和需要恢复的区域。我们可以直接拿走该区域，也可以留下位置而改变其中的数值；模型分别面对内容缺失和内容受损两种输入。前一种构造称为遮蔽恢复，后一种构造用于去噪恢复。把选中内容替换成缺失符号也是一种扰动，因此这里分开比较的是常用输入安排及损失位置，并非两个互斥的理论类别。

沿图~\ref{fig:data-derived-knowledge-masking-denoising}的同一$3\times3$网格观察这两种输入：上行两处原值被叉号替代，下行相同位置仍有数值，但已被改动。两行的保留目标都没有随输入修改。赭色边框指定本例只核对这两处输出；若实际去噪任务核对全部位置，未受损内容也会参与梯度计算，目标定义必须反映这一变化。

\begin{figure*}[htbp]
  \centering
  \begin{tikzpicture}[
  x=1mm,y=1mm,
  font=\figurefont\scriptsize,
  flow/.style={knowledge arrow,draw=NNWire,line width=1.1pt},
  grid/.style={draw=BookBlue,line width=.85pt},
  cell/.style={draw=BookRule,line width=.5pt,minimum size=5mm,inner sep=0pt},
  heading/.style={font=\figurefont\bfseries\footnotesize,text=BookInk},
  stage/.style={font=\figurefont\footnotesize,text=BookMuted}
]
  \node[heading,anchor=east] at (13,1) {掩码};
  \node[heading,anchor=east] at (13,-34) {去噪};
  \foreach \x/\name in {23/原始观测,64/处理后预测输入,105/模型预测,146/保留目标}{
    \node[stage] at (\x,13) {\name};
  }

  \begin{scope}[shift={(15,8)}]
    \foreach \r/\c/\v in {0/0/.2,0/1/.4,0/2/.6,1/0/.3,1/1/.8,1/2/.7,2/0/.1,2/1/.5,2/2/.9}{
      \node[cell] at ({\c*5},{-\r*5}) {$\v$};
    }
  \end{scope}
  \begin{scope}[shift={(56,8)}]
    \foreach \r/\c/\v in {0/0/.2,0/1/.4,0/2/{\times},1/0/.3,1/1/{\times},1/2/.7,2/0/.1,2/1/.5,2/2/.9}{
      \node[cell] at ({\c*5},{-\r*5}) {$\v$};
    }
    \fill[FigureOutput!18] (7.5,-2.5) rectangle (12.5,2.5);
    \fill[FigureOutput!18] (2.5,-7.5) rectangle (7.5,-2.5);
    \node[cell] at (10,0) {$\times$};
    \node[cell] at (5,-5) {$\times$};
  \end{scope}
  \begin{scope}[shift={(97,8)}]
    \foreach \r/\c/\v in {0/0/.2,0/1/.4,0/2/.55,1/0/.3,1/1/.70,1/2/.7,2/0/.1,2/1/.5,2/2/.9}{
      \node[cell] at ({\c*5},{-\r*5}) {$\v$};
    }
    \draw[FigureLoss,line width=1.1pt] (7.5,-2.5) rectangle (12.5,2.5);
    \draw[FigureLoss,line width=1.1pt] (2.5,-7.5) rectangle (7.5,-2.5);
  \end{scope}
  \begin{scope}[shift={(138,8)}]
    \foreach \r/\c/\v in {0/0/.2,0/1/.4,0/2/.6,1/0/.3,1/1/.8,1/2/.7,2/0/.1,2/1/.5,2/2/.9}{
      \node[cell] at ({\c*5},{-\r*5}) {$\v$};
    }
    \draw[FigureLoss,line width=1.1pt] (7.5,-2.5) rectangle (12.5,2.5);
    \draw[FigureLoss,line width=1.1pt] (2.5,-7.5) rectangle (7.5,-2.5);
  \end{scope}
  \draw[flow] (31,3) -- (53,3);
  \draw[flow] (72,3) -- (94,3);
  \draw[flow] (113,3) -- node[above,text=FigureLoss] {仅比较框内位置} (135,3);

  \begin{scope}[shift={(15,-27)}]
    \foreach \r/\c/\v in {0/0/.2,0/1/.4,0/2/.6,1/0/.3,1/1/.8,1/2/.7,2/0/.1,2/1/.5,2/2/.9}{
      \node[cell] at ({\c*5},{-\r*5}) {$\v$};
    }
  \end{scope}
  \begin{scope}[shift={(56,-27)}]
    \foreach \r/\c/\v in {0/0/.2,0/1/.4,0/2/.8,1/0/.3,1/1/.6,1/2/.7,2/0/.1,2/1/.5,2/2/.9}{
      \node[cell] at ({\c*5},{-\r*5}) {$\v$};
    }
    \fill[BookAmber!18] (7.5,-2.5) rectangle (12.5,2.5);
    \fill[BookAmber!18] (2.5,-7.5) rectangle (7.5,-2.5);
    \node[cell] at (10,0) {.8};
    \node[cell] at (5,-5) {.6};
  \end{scope}
  \begin{scope}[shift={(97,-27)}]
    \foreach \r/\c/\v in {0/0/.2,0/1/.4,0/2/.65,1/0/.3,1/1/.72,1/2/.7,2/0/.1,2/1/.5,2/2/.9}{
      \node[cell] at ({\c*5},{-\r*5}) {$\v$};
    }
    \draw[FigureLoss,line width=1.1pt] (7.5,-2.5) rectangle (12.5,2.5);
    \draw[FigureLoss,line width=1.1pt] (2.5,-7.5) rectangle (7.5,-2.5);
  \end{scope}
  \begin{scope}[shift={(138,-27)}]
    \foreach \r/\c/\v in {0/0/.2,0/1/.4,0/2/.6,1/0/.3,1/1/.8,1/2/.7,2/0/.1,2/1/.5,2/2/.9}{
      \node[cell] at ({\c*5},{-\r*5}) {$\v$};
    }
    \draw[FigureLoss,line width=1.1pt] (7.5,-2.5) rectangle (12.5,2.5);
    \draw[FigureLoss,line width=1.1pt] (2.5,-7.5) rectangle (7.5,-2.5);
  \end{scope}
  \draw[flow] (31,-32) -- (53,-32);
  \draw[flow] (72,-32) -- (94,-32);
  \draw[flow] (113,-32) -- node[above,text=FigureLoss] {此例比较受扰位置} (135,-32);

  \node[font=\figurefont\scriptsize,text=BookMuted,anchor=west] at (15,-48) {数值与扰动均为教学构造；保留目标是原始观测，不代表无采集噪声的真值。};
\end{tikzpicture}

  \caption{在同一网格上对照隐藏内容与修改内容。叉号标出掩码，浅赭色标出受扰位置，赭色边框指定本例计分范围；两行都用未修改的原记录核对预测。数值与扰动仅为教学构造，原记录还可能含有采集噪声，不代表无噪声真值。}
  \label{fig:data-derived-knowledge-masking-denoising}
\end{figure*}
\FloatBarrier

\subsubsection{基于掩码的内容恢复}
\label{sec:data-derived-knowledge-masked-recovery}

在图中上行，我们要让模型补出两个已经拿走的值，而不让它直接复制原值。\term{掩码恢复}（masked reconstruction）就是先选定待恢复位置$M$，再只根据其余位置$V=I\setminus M$预测$\bm{x}_M$。隐藏单位可以是文本词元、不重叠图像块、表格字段或单元格。单位大小改变了可用线索：单个像素常可由邻域插值补出，连续区域或完整属性则可能需要较远的上下文。

要根据这些保留值更新模型，还需把每处恢复偏差汇成一个可计算的分数。对非空掩码集合$M$中的连续像素或像素块，用平方欧氏距离度量预测与原值的偏差，再对选中单位取平均；训练在所用网络结构允许的参数范围内减小这一分数。记$\bm{x}_j$为目标单位的像素向量、$\widehat{\bm{x}}_j$为对应预测，常见的掩码均方误差为
\begin{equation}
  \mathcal{L}_{\mathrm{mask}}
  =\frac{1}{|M|}\sum_{j\in M}\lVert \widehat{\bm{x}}_j-\bm{x}_j\rVert_2^2.
  \label{eq:data-derived-knowledge-masked-reconstruction}
\end{equation}
式~\eqref{eq:data-derived-knowledge-masked-reconstruction}越小，选中位置的平均恢复偏差越小；分母$|M|$使计分不直接随隐藏单位数增长。离散词元没有同样的像素距离，可以改用真实词元的交叉熵计分。若把未遮蔽位置也纳入平均，模型复制大量可见内容便能降低总分。这种目标可以使用，但它同时奖励复现已知内容，不能再把总分的降低全部解释为恢复未知内容的进步。

现在用图~\ref{fig:data-derived-knowledge-masking-denoising}上行复算这一计分。两个保留值为$(0.6,0.8)$，人为设定的预测为$(0.55,0.70)$；把它们代入式~\eqref{eq:data-derived-knowledge-masked-reconstruction}，得到
\[
  \mathcal{L}_{\mathrm{mask}}
  =\frac{(0.55-0.6)^2+(0.70-0.8)^2}{2}
  =0.00625.
\]
得到的$0.00625$只核对两处隐藏值，其余七处可见值没有参与平均。这里恢复的是训练记录中的$0.6$和$0.8$，不是经额外测量确认的无噪声真值；这正是观测目标所能提供的答案边界。

把这一构造用于图像编码时，掩码自编码器（masked autoencoder, MAE）让编码器只处理可见块，再由轻量解码器结合掩码标记补出缺失块的像素。原论文只在被遮蔽块上计算均方误差\scite{ka-he2022mae}
编码器读取的块与损失核对的块因而互补。按块归一化后的像素也可作为目标，MAE对此作过研究；归一化改变数值尺度，却没有改变答案来自原图这一事实。

可见块变少以后，恢复任务会改变，而不只是多算几个误差项。遮蔽比例过低可能只要求局部插值；比例过高则可能让同一可见输入支持多个合理答案，平方误差下的点预测倾向平均这些答案。随机、连续片段和结构化遮蔽又各自留下不同上下文。因此需要同时考察比例与位置规则：让恢复更容易虽然通常能降低训练损失，却未必让编码器保留更多物种识别信息。

仅在最终输入上画出掩码，还不足以隔离答案。若此前的平滑、插值或全图归一化统计读过目标位置，邻居数值或统计量可能已携带被隐藏的内容。我们需要逐个核对预处理算子读取哪些位置，并在输入与目标拆分后检查它是否仍访问$M$；遮蔽标记本身不能消除较早步骤造成的泄露。

输入仍可使用掩码，而答案也可以改成另一编码器算出的特征。此时改变的是目标来源：需要确认该编码器是否固定、是否随训练更新、是否接收当前损失的梯度，不能再从“做了遮蔽”推出“预测的是观测原值”。其中随训练变化的表示目标由第~\ref{sec:data-derived-knowledge-training-objectives}节接着分析。

\subsubsection{基于扰动的去噪恢复}
\label{sec:data-derived-knowledge-denoising-recovery}

图中下行没有拿走位置，而是把两处原值改成了受损数值。我们仍用未改动的原记录核对输出，这称为\term{去噪恢复}（denoising reconstruction）：给定$\bm{x}$，按扰动分布$q(\widetilde{\bm{x}}\mid\bm{x})$生成输入$\widetilde{\bm{x}}$，再预测扰动前的内容。扰动可采用加性噪声、词元替换、局部打乱或颜色缺失等人为损坏；这里要撤销的是构造步骤的损坏，不是宣称已经测得采集过程的物理噪声。

同一原记录可生成多种损坏版本，一次恢复误差不足以评价模型对整套扰动规则的适应。我们先对选定位置的偏差取平均，再对$q$允许生成的输入取期望，使训练减小这些损坏版本上的平均误差，而非只拟合一个副本。设非空$R\subseteq I$为已知受扰位置，$f_{\bm{\theta}}$为恢复模型，$\ell$按目标类型计分；固定原记录$\bm{x}$时，一般形式为
\begin{equation}
  \mathcal{L}_{\mathrm{denoise}}
  =\mathbb{E}_{\widetilde{\bm{x}}\sim q(\cdot\mid\bm{x})}
   \left[\frac{1}{|R|}\sum_{j\in R}\ell\!\left(f_{\bm{\theta}}(\widetilde{\bm{x}})_j,x_j\right)\right].
  \label{eq:data-derived-knowledge-denoising}
\end{equation}
式~\eqref{eq:data-derived-knowledge-denoising}在$R$上计分，回答受扰位置恢复得怎样。若把平均范围改为整个$I$，模型还要保持未扰内容，梯度的组成也随之改变。报告去噪结果时，因而需要给出扰动分布和损失范围，单写一个损失名称不能确定实际学了什么。

恢复误差还受输入剩余信息限制。扰动很弱时，接近恒等复制的策略可能已经使误差很小；一块花瓣完全被随机纹理替换后，剩余输入却可能无法确定其颜色。后一种情况下，模型按所用损失输出条件均值、众数或一个分布，不能仅靠更充分的优化找回未被输入确定的答案。均方误差偏好多个合理颜色的平均，分布预测则可保留多种可能。

要反复抽取损坏版本，可以每轮在线生成，也可以预先保存多个副本。前者保存原样本和随机配置即可，但每轮增加处理；后者减少重复生成，却占用存储并把可用扰动限制在已保存副本内。随机种子或分布参数应随构造记录保留，评价也需说明使用相同规则还是另一种扰动。在同一种教学噪声上恢复准确，只支持该规则下的结果，不能推出能够修复真实采集误差。

即使恢复误差下降，也还没有回答花卉识别是否受益。准确撤销固定压缩伪影的模型可能没有保留花瓣形态；像MAE这类编码器的像素恢复也可能不唯一，却仍能产生用于识别的特征。恢复质量评价核对原记录，下游评价核对任务所需信息，两种证据必须分开取得。

\subsection{基于上下文预测的目标构造}
\label{sec:data-derived-knowledge-context-prediction}

对记录“花朵逐渐开放”，我们可以给模型前缀“花朵逐渐”，把“开放”留下作答案。记录已有的顺序为切分提供了依据，无需随机挑选隐藏位置。\term{上下文预测}（context prediction）指根据约定前缀、邻域或其他可见内容预测目标片段。它与掩码恢复可能都预测词元，但允许读取哪一侧内容、怎样切出目标，是两个不同的选择。

为了评价各个前缀是否支持记录中的下一词元，需要模型给出候选词元的条件概率。对序列$x_1,\ldots,x_T$，第$t$处只读取$x_{<t}$，在所用词表中为下一词元分配概率；真实词元概率越高，负对数计分越低。对全部约定位置求和，便把一条记录的预测要求汇成逐位置负对数似然；$t=1$时前缀为空或由约定的起始标记表示：
\begin{equation}
  \mathcal{L}_{\mathrm{context}}
  =-\sum_{t=1}^{T}\log p_{\bm{\theta}}(x_t\mid x_{<t}).
  \label{eq:data-derived-knowledge-context-prediction}
\end{equation}
式~\eqref{eq:data-derived-knowledge-context-prediction}把每个真实词元接回其可见前缀，与\BookSectionRef{sec:pgm-static-directed-autoregressive}{卷6《模型：概率图模型》第11章“静态有向模型”第11.1.2节“自回归密度模型”}的自回归分解一致。就训练依据而言，“花朵逐渐开放”可提供“花朵$\to$逐渐”和“花朵逐渐$\to$开放”等位置目标，答案都来自原记录。增加位置扩展了监督覆盖，也增加了相应预测计算，计算量近似随位置数增长。

如果改为遮住句中的一个词元，模型通常还能读取其左右两侧；前缀预测只允许使用之前的内容。BERT的掩码词元目标采用前一种安排，以双向上下文预测从记录中保留的原词元\scite{devlin2019}
两种目标的答案可能相同，可见条件却不同，使用交叉熵不会使两个条件概率互换。

式中训练输入通常是真实前缀$x_{<t}$，这种安排称为教师强制（teacher forcing），即每一步都用记录中的前文而非模型先前生成的词元。连续生成时，后续位置却可能接到已有预测，前一步错误也会改变后一步输入。因此训练时的一步误差不能直接代替长序列生成误差。截断窗口还会移除较早信息，位置和片段边界的选择也可能让训练上下文偏离真实使用情形。

从文本推广到其他记录时，仍可沿用“限制可见内容、保留记录片段”的构造。图像可用左侧区域预测右侧区域，音频可用邻近频段预测保留片段；一个模态也可预测同一事件中已经记录的另一模态。变化的是输入与目标的数据类型，保留下来的判断依据是答案能否直接从记录取出。若预测的是另一模态编码后的向量，则还需追踪编码器，不再仅凭原始配对判断目标来源。

保留相邻内容不保证模型必须利用主体信息。网页模板、固定边框和周期背景可能已经足以完成预测；分词器、缓存特征或裁剪坐标若携带答案，也会让训练计分偏低。改变上下文范围或删除可疑元数据可检查模型依赖什么，再用独立下游任务检查这些规律是否有用，才能区分复现记录规律与取得预期任务知识。

\subsection{基于时序预测的目标构造}
\label{sec:data-derived-knowledge-temporal-prediction}

假设我们已有一段植物生长影像，要用今天及此前的帧预测几天后的状态。训练数据整理时可以拿到后续帧，真正站在今天作出判断时却不能读取它们。\term{时序预测}（temporal prediction）就是以真实时点$t$为信息边界，用截至$t$的历史预测之后的记录。与仅由文本排列决定的上下文顺序相比，这里额外约束信息必须在判断时点可获得；未来内容虽仍是观测目标，却只能供训练核对。

“预测未来”还没有指定哪一段历史与哪一段未来相配，无法据此切出训练样本或比较误差。我们需要用窗口明确这个对应：设$L\geq1$为历史窗口长度，$h\geq1$为起始预测跨度，$H\geq1$为目标轨迹长度，均以采样步计。冒号表示包含两端的连续索引，输入与保留目标写为
\begin{equation}
  \bm{x}_{t-L+1:t}
  \quad\longrightarrow\quad
  \bm{y}_{t,h,H}
  =\bm{x}_{t+h:t+h+H-1}.
  \label{eq:data-derived-knowledge-temporal-windows}
\end{equation}
式~\eqref{eq:data-derived-knowledge-temporal-windows}固定了允许读取的最后时点和目标起点，但步数还需配合采样间隔解释：每分钟和每天采样的“未来十步”不是同一个问题。$H=1$只核对一个时点，$H>1$则核对一段植物生长轨迹，输出维数增大，多个位置的误差也可能相关。

确定窗口后，还需决定怎样表达未来答案。当未来近似唯一，且平方误差符合实际使用代价时，点预测是可用选择；同一历史可能通向多条轨迹时，条件分布、分位数或离散区间能保留更多不确定性。分布输出没有使未来变得可预测，只是没有把全部可能压成均值，评价规则仍须对应实际决策。

把$h$增大，模型要跨过更长的未观测过程。短跨度常可利用局部连续性，长跨度更容易受到干预、累积误差和环境变化影响。比较时应固定数据范围，分别报告各跨度的误差，避免一个平均值掩盖远期失效；同时训练多个跨度可扩展目标覆盖，但要承担更多输出计算和目标存储。

判断时点的边界必须贯穿切分和预处理。若先随机打散窗口，同一轨迹的高度重叠片段可能同时进入训练和测试；若用全序列均值归一化，或用未来点插值当前缺失值，输入也会包含未来信息。应先按照预期使用单位划分时间、植株、设备或轨迹，再只用训练阶段允许获得的统计量拟合预处理，并检查每个输入在其判断时点是否确实可计算。

在这个边界内，较低预测误差仍可能来自重复背景，而非植物状态；长期过程本身随机性较大时，较高误差也不能证明编码毫无价值。我们需要分开回答未来能被现有输入确定到什么程度、模型实际保留了什么信息，以及这些信息是否支持下游用途，而不从单个预测分数直接认定知识质量。

\section{基于样本关系的知识获取}
\label{sec:data-derived-knowledge-relations}

现在仍用这批花卉照片，但不再要求补出花瓣像素。若两个裁剪都保留同一朵花，我们可以要求模型保留它们的共同信息；若还需要识别哪一个裁剪与当前照片对应，就要提供其他候选作参照。\term{关系知识}（relational knowledge）是这种可追溯的关系记录：正对应说明希望共享什么，负对照说明当前判别允许区分谁。关系可由变换、采集记录或采样规则取得并在编码前保存，参加比较的表示数值却要等当步模型计算后才出现。

正对应还没有唯一规定训练目标。对于同一对裁剪，我们可以在候选中提高配对的相似度，也可以预测另一侧向量，或同时改变两路使它们接近，再限制一批表示的变化。三种要求使用同一关系，却比较不同的量并产生不同梯度。反过来，损失中出现余弦相似度，也不能说明裁剪的关系依据是否可靠。

因此，我们先建立关系，再把关系接到模型输出，最后检查相邻训练步怎样改变目标与参数。可保存的关系记录、当步算出的目标值和跨步更新的参数是三个对象：构造证据检查关系假设，梯度路径检查目标是否按预期使用关系，下游评价检查学得的差异是否有用。一次损失下降不足以同时回答这三个问题。

\subsection{样本关系的构造}
\label{sec:data-derived-knowledge-relation-construction}

要把“两个裁剪属于一对”交给训练，至少要知道原对象、两侧记录、配对或排除依据、适用范围和构造配置。同源裁剪还需检查是否都覆盖主体，同一时刻的音频与画面还需规定同步误差。来源标识只能帮助追溯，不能替代这些成立条件；暂时无法确定的关系可以保留置信范围，而不必强行记成是或否。

图~\ref{fig:data-derived-knowledge-relation-construction}把构造动作与关系判断分开画出。蓝色箭头只表示生成视图，青色实线才表示预期共享信息的正对应；左侧灰色虚线指定待检查的负对照，右侧赭色虚线指定教学例子的错配。沿线型区别，我们可以分别问裁剪是否保留主体、候选是否允许被排斥、描述是否覆盖照片，而不从程序生成或采集配对直接推出语义正确。

\begin{figure*}[htbp]
  \centering
  \begin{tikzpicture}[
  x=1mm,y=1mm,
  font=\figurefont\footnotesize,
  panel/.style={draw=BookRule,fill=NNPanel,line width=.85pt,rounded corners=1.5mm},
  item/.style={draw=FigureInput,fill=FigureInput!10,line width=1.1pt,text width=25mm,minimum height=11mm,align=center,inner sep=1.5mm},
  textitem/.style={draw=FigureOutput,fill=FigureOutput!8,line width=1.1pt,text width=28mm,minimum height=11mm,align=center,inner sep=1.5mm},
  generate/.style={knowledge arrow,draw=FigureInput,line width=1.1pt,rounded corners=1mm},
  positive/.style={draw=FigureModel,line width=1.4pt},
  negative/.style={draw=BookMuted,line width=.9pt,dashed,rounded corners=1mm},
  warning/.style={font=\figurefont\scriptsize,text=FigureLoss,align=center}
]
  \draw[panel] (1,-77) rectangle (82,12);
  \draw[panel] (87,-77) rectangle (168,12);
  \node[font=\figurefont\bfseries\footnotesize,text=BookInk] at (41.5,5) {(a) 观测变换建立正对应};
  \node[item] (original) at (20,-13) {原照片\\花朵与叶片};
  \node[item] (crop-a) at (61,-13) {裁剪A\\花朵};
  \node[item] (crop-b) at (61,-37) {裁剪B\\花朵与叶片};
  \node[item,draw=BookAmber,fill=BookAmber!10] (same-class) at (20,-63) {同类另一株\\未必是负语义};
  \coordinate (crop-split) at (40,-13);
  \draw[FigureInput,line width=1.1pt] (original.east) -- (crop-split);
  \draw[generate] (crop-split) -- (crop-a.west);
  \draw[generate] (crop-split) |- (crop-b.west);
  \draw[positive] (crop-a.south) -- node[right=1mm,text=FigureModel] {正对应} (crop-b.north);
  \draw[negative] (crop-b.south) |- (same-class.east);
  \node[font=\figurefont\scriptsize,text=BookMuted,anchor=south] at (49,-61) {候选负对照};
  \node[warning,text width=28mm] at (20,-37) {同源裁剪也可能\\丢失共同主体};

  \node[font=\figurefont\bfseries\footnotesize,text=BookInk] at (127.25,5) {(b) 采集日志建立自然配对};
  \node[item] (photo) at (104,-13) {花卉照片\\白色花瓣};
  \node[textitem] (caption) at (148,-13) {原有描述\\“白色花瓣，长叶”};
  \node[textitem,draw=BookAmber,fill=BookAmber!10] (wrong-caption) at (148,-41) {错配描述\\“黄色花瓣”};
  \draw[positive] (photo.east) -- (caption.west);
  \node[font=\figurefont\scriptsize,text=FigureModel,anchor=south] at (125,-11) {观测配对};
  \draw[negative,draw=FigureLoss] (photo.south) |- (wrong-caption.west);
  \node[warning,anchor=south] at (118,-39) {错误配对};
  \node[warning,text width=70mm] at (127.5,-64) {原有描述可能只覆盖部分内容；\\错配样本不能作为正对应。};

  \draw[generate] (3,-84) -- (16,-84);
  \node[anchor=west,text=BookMuted] at (18,-84) {生成视图};
  \draw[positive] (53,-84) -- (66,-84);
  \node[anchor=west,text=BookMuted] at (68,-84) {正对应};
  \draw[negative] (103,-84) -- (116,-84);
  \node[anchor=west,text=BookMuted] at (118,-84) {负对照或错配关系};
\end{tikzpicture}

  \caption{生成视图与确认关系是不同动作。左侧同源裁剪形成候选正对应，同类另一株还需判断是否允许作负对照；右侧已有图文配对与错配描述形成对照。连为负对照没有确认语义不同，已知错配则不能拿来提供正对应。}
  \label{fig:data-derived-knowledge-relation-construction}
\end{figure*}
\FloatBarrier

\subsubsection{基于数据变换的正对应关系构造}
\label{sec:data-derived-knowledge-transformation-positive}

同一张花卉照片可随机裁剪两次，从而不经逐对人工标注便得到一对视图。这里的依据是变换保留希望共享的内容，同时改变不希望表示依赖的因素。整体与局部视图、尺度变化、颜色扰动，以及其他数据中的语音加噪或文本改写，改变的内容并不相同，需要分别确定保留条件和强度。

两次裁剪都含有花朵时，匹配可鼓励模型忽略取景差异；一路只剩背景时，却未必还有足够共同内容。颜色扰动若消除了品种判断需要的花色，要求表示保持不变就与该任务冲突。同样，方向或位置属于下游答案的一部分时，旋转、翻转和大范围平移也不能无条件作为不改变表示的变换。

这些反例说明，变换配置实际声明了哪些变化不应影响表示。记录种类、组合顺序、强度、随机参数和主体覆盖情况，才能用抽样可视化或自动规则检查共同内容是否仍在。增加视图会增加对应关系，也会增加编码次数及相应计算，不能只计节省的人工关系标注。

若把一对视图扩展为一个整体视图与多个局部视图，关系范围也随之改变。每个局部只覆盖整体的一部分，“来自同一图像”可以继承，却不能据此要求局部完整复制整体信息。局部到全局的匹配是否可用，还依赖多裁剪设计与训练经验，来源一致本身没有提供成功保证。

到这里取得的是正对应及其范围，尚没有规定用候选判别还是直接匹配，也没有指定独立目标分支。第~\ref{sec:data-derived-knowledge-training-objectives}节会把这些已确定的关系接到具体损失与更新路径。

\subsubsection{基于自然配对的正对应关系构造}
\label{sec:data-derived-knowledge-natural-positive}

照片若已带有说明文字，我们可以复用这条关联，而不必从同一像素记录生成两侧。\term{自然配对}（natural pairing）指采集过程中已有的对应，如同一对象的多次拍摄、同步音视频、配对图文或相邻时刻读数。两侧可以属于不同数据类型，配对依据则从变换来源转为对象标识、时间戳、事件编号或文档结构，需要检查这些记录是否可靠。

同一张花卉照片的文字可能描述整株形态，也可能只描述花色，还可能误配自另一照片。前两种配对分别支持较广与较窄的共同内容，第三种是关系错误。因此训练前要明确整例、局部或一对多的对应范围，并保留缺失一侧及多重匹配，不能把“有配对记录”统一解释成内容逐项相同。

% 完整作者和会议条目不能跨页续排。
\Needspace{65mm}
CLIP把已有图文配对作为正对应，再以同批次其他图文组合作参照，训练模型区分哪个组合与记录中的配对相符\scite{radford2021}
这复用了人类书写的描述信息，并非重新为目标分类逐图标注。节省后一种投入不等于完全没有人类监督，描述来源、关联核对和任务标注的成本也不能混为一项。

配对错误也会进入后续判别。多张照片复用同一模板描述时，彼此作为负对照可能产生冲突；音频落后于画面时，相同时刻的记录可能来自不同事件。时间偏移、描述不完整、重复和共享背景都需要回查采集协议与关联记录，增加同样方式采集的数据不会自动纠正系统性错配。

两侧数据还要分别编码和存储。报告各侧缺失率、配对核验方式与模态处理开销，才能说明复用已有关联的实际投入；“无需新标签”只回答其中一种成本是否增加。

\subsubsection{基于样本采样的负对照关系构造}
\label{sec:data-derived-knowledge-negative-sampling}

若要在多个裁剪中找到当前照片的对应裁剪，只有一条正对应还不够，必须确定允许比较哪些其他候选。从小批量、缓存或检索集合选出的\term{负对照}（negative contrast）承担这种参照角色：当前目标允许查询相对它们取得更高匹配分数。被选入集合，不等于已经得到语义不同的人工判断。

随机抽取可覆盖较广候选，但差异很大的记录可能只产生较弱的区分梯度。按当前表示检索相似记录，可得到\term{难负例}（hard negative），即模型此时难以与查询区分的候选，训练便更多关注细微差别。难度来自当前表示，不证明语义不同；反复检索少数记录还可能强化这套表示已有的偏差。

回到花卉照片，同一植株的另一张照片、同品种另一株与完全不同对象不能固定承担同一种角色。识别具体植株时，前两者分别可能是正对应和可区分实例；识别品种时，排斥同品种样本却可能违背预期聚合。应先确定关系服务的任务和范围，“尚未记录为正例”不能直接成为“确定负例”的依据。

即使没有单独保存采样文件，小批量组成也已经选择了进入对比目标的候选。重复、类别失衡和同源记录会改变参照分布，缓存中的旧参数表示则可能偏离当前模型。扩大集合增加了区分要求，也增加错误负例机会、通信和存储，集合越大并非无条件越有用。

一个查询若有多个合理对应，可以将它们显式记为正对应，或采用支持多正例的目标。否则遗漏的正对应会作为其他候选进入分母而受到排斥。关系集合决定训练允许区分谁，损失决定怎样施加区分；前者的错误不能仅靠后者下降来排除。

\subsection{训练目标的构造}
\label{sec:data-derived-knowledge-training-objectives}

已有两路正对应输入后，我们还缺少用于参数更新的计分规则。可以比较它们相对其他候选的相似度，预测另一侧输出，或匹配两路并检查一批表示的统计量。为确定训练实际改变什么，以下逐条追踪输入、附加映射、比较量和梯度，再把训练头与下游特征分开。尤其需要区分交换视图另算一项，与同一项损失同时向两路回传，前者不自动包含后者。

% 标题、表示定义及其边注在同一页内保留足够高度。
\Needspace{80mm}
\subsubsection{基于对比判别的目标构造}
\label{sec:data-derived-knowledge-contrastive}

假设当前查询是照片$A$的一个裁剪，候选包含它的另一裁剪和其他照片的视图，我们要让已知配对取得相对优势。\term{对比判别}（contrastive discrimination）把这个要求写成正对应相对候选具有更高相似度的目标。SimCLR让同批$n$个原样本各生成两个视图，得到$2n$个编码：$f_{\bm{\theta}}$产生特征$\bm{h}_i$，投影器$g_{\bm{\theta}}$再把它映射到专供训练计分的$\bm{z}_i=g_{\bm{\theta}}(\bm{h}_i)$。训练后通常保留$\bm{h}_i$作下游特征，因此改善投影空间的配对分数与使用编码特征是不同环节\scite{chen2020}

要比较配对优势，先需给任意两个投影向量一个同口径分数。这里用余弦相似度比较方向，不把向量长度直接作为优势；以下假定$n\geq2$，且参与比较的向量非零，使分母有定义：
\[
  \operatorname{sim}(\bm{z}_i,\bm{z}_k)
  =\frac{\langle\bm{z}_i,\bm{z}_k\rangle}
  {\lVert\bm{z}_i\rVert_2\lVert\bm{z}_k\rVert_2}.
\]
再固定一个锚点，即当前被查询的视图$i$。它的候选由本批其余$2n-1$个视图组成，排除自身并包含唯一正对应$j$；我们用指数化分数在这组候选上归一化，再惩罚正对应的低概率。温度$\tau>0$控制相似度差异在归一化中被放大的程度，负对数使正确配对概率越大、损失越小。SimCLR的单向归一化温度交叉熵损失为
\begin{equation}
  \ell_{i,j}
  =-\log
  \frac{\exp\!\left(\operatorname{sim}(\bm{z}_i,\bm{z}_j)/\tau\right)}
  {\displaystyle\sum_{k=1}^{2n}\mathbb{I}[k\neq i]\exp\!\left(\operatorname{sim}(\bm{z}_i,\bm{z}_k)/\tau\right)},
  \qquad \tau>0.
  \label{eq:data-derived-knowledge-simclr}
\end{equation}
式~\eqref{eq:data-derived-knowledge-simclr}中$\mathbb{I}[k\neq i]$只排除自身，正对应$j$与其余$2n-2$个对照都参与归一化。完整目标对$2n$个视图轮流作锚点的损失取平均，梯度回到共享编码器与投影器。减小温度会强调相似度差异，也会增强难负例和错误关系的影响，不能把低温仅理解为更精确。

用两个原样本可以直接检查候选范围。教学批次取$n=2$，按$A_1,A_2,B_1,B_2$排列四个视图，配对为$(A_1,A_2)$与$(B_1,B_2)$；对角线是自身比较，配对位置是正对应，其余位置在本例中作负对照。人为设定余弦相似度矩阵为
\[
  \bm{S}=
  \begin{bmatrix}
  1&0.8&0.1&0\\
  0.8&1&0&-0.1\\
  0.1&0&1&0.7\\
  0&-0.1&0.7&1
  \end{bmatrix}.
\]
取$\tau=0.5$，锚点$A_1$只与$A_2,B_1,B_2$比较，故配对概率为$\exp(1.6)/[\exp(1.6)+\exp(0.2)+\exp(0)]\approx0.690$，负对数损失为$\ell_{A_1,A_2}\approx0.371$。这个概率回答当前三名候选中的配对支持程度，没有定义花卉类别，因而不是同一类别的概率。

这一算例也说明，小批量改变后，即使同一正对应的相似度不变，分母和损失仍会改变。更大批量通常需要更多相似度计算，分布式训练还需交换表示；候选缺少区分力时，小批量目标可能较弱。缓存或检索能扩大参照范围，但要额外检查旧表示与采样偏差。

若把同品种花卉误作不应匹配的候选，它仍会进入式~\eqref{eq:data-derived-knowledge-simclr}的分母。实例区分于是可能排斥物种识别希望聚合的样本。多正例目标、去重与关系过滤可减轻这种冲突，却都依赖额外关系信息，损失本身不知道哪一种语义应被保留。

还可以用一个极端反例检查目标是否真的奖励样本差异。若所有投影都等于同一非零向量$\bm{u}$，每项余弦相似度为$1$，候选无法区分，代入损失得到
\[
  \ell_{i,j}=\log(2n-1).
\]
常数输出没有配对优势。当模型及表示空间允许正对应更相似时，可以取得更小损失，上述单项就有$0.371<\log3\approx1.099$。但可取得较低值不等于优化会离开常数状态：相同非零向量的余弦相似度对任一输入的梯度为零，精确常数输出仍是关于表示的驻点，也就是当前一阶梯度不能推动表示变化的位置。若$n=1$，唯一候选就是正对应，损失更恒为零。对比目标因此只在有效候选和可实现区分的条件下偏好差异，既不保证任意初始化都避免坍塌，也不保证差异来自花朵而非背景、设备或水印。

\subsubsection{基于跨视图预测的目标构造}
\label{sec:data-derived-knowledge-cross-view-prediction}

仍给定同一照片的两个裁剪，现在不要求从其他候选中找配对，而让裁剪$A$预测裁剪$B$的特征。这是\term{跨视图预测}（cross-view prediction）：目标由对应视图的编码分支计算，只使用正对应。BYOL用在线分支执行编码$f_{\bm{\theta}}$、投影$g_{\bm{\theta}}$和预测$q_{\bm{\theta}}$，目标分支则只有参数为$\bm{\xi}$的编码器与投影器。附加预测器把在线投影映射为待匹配输出；对$\bm{v}_A,\bm{v}_B$分别记为
\[
  \bm{p}_A=q_{\bm{\theta}}\!\left(g_{\bm{\theta}}(f_{\bm{\theta}}(\bm{v}_A))\right),
  \qquad
  \bm{z}^{\mathrm{tar}}_B=g_{\bm{\xi}}(f_{\bm{\xi}}(\bm{v}_B)).
\]
这两路向量需要同一个偏差分数，才能让在线网络按目标改进预测。假定$\bm{p}_A$与$\bm{z}^{\mathrm{tar}}_B$维数相同且均非零，先作$\ell_2$归一化以比较方向，再用平方距离计分；距离越小，当前预测与目标越一致。$\operatorname{sg}$表示停止梯度（stop-gradient），即前向保留目标数值，当前反向却将它视作常量。单向损失为
\begin{equation}
  \ell_{A\to B}
  =\left\lVert
  \frac{\bm{p}_A}{\lVert\bm{p}_A\rVert_2}
  -\operatorname{sg}\!\left(
  \frac{\bm{z}^{\mathrm{tar}}_B}{\lVert\bm{z}^{\mathrm{tar}}_B\rVert_2}
  \right)\right\rVert_2^2
  =2-2\frac{\langle\bm{p}_A,\bm{z}^{\mathrm{tar}}_B\rangle}
  {\lVert\bm{p}_A\rVert_2\lVert\bm{z}^{\mathrm{tar}}_B\rVert_2}.
  \label{eq:data-derived-knowledge-byol}
\end{equation}
等式右侧由两个单位向量的平方范数均为$1$得到，说明减小距离等价于提高它们的余弦相似度。BYOL交换输入再计算$\ell_{B\to A}$，把两项相加；每项仍只向相应在线分支求导。两向求和让两个视图都承担预测输入，却没有把目标端改成接受梯度的一端\scite{ka-grill2020byol}

目标端不接受当前损失梯度，仍需说明下一步的目标从哪套参数产生。BYOL不用该梯度直接更新目标，而在在线更新完成后，对目标编码器、投影器与在线对应参数作指数移动平均（exponential moving average, EMA），即保留一部分旧值并混入一部分新值。迭代步$t$的更新为
\begin{equation}
  \bm{\xi}^{(t+1)}
  =m\bm{\xi}^{(t)}+(1-m)\bm{\theta}_{\mathrm{enc+proj}}^{(t+1)},
  \qquad 0\leq m\leq1.
  \label{eq:data-derived-knowledge-byol-ema}
\end{equation}
式~\eqref{eq:data-derived-knowledge-byol-ema}中的$\bm{\theta}_{\mathrm{enc+proj}}^{(t+1)}$只包含更新后的编码器与投影器，预测器没有目标侧对应参数，只作梯度更新。停止梯度使当步答案不被这项误差推动，EMA则令以后生成答案的参数逐步跟随在线网络；两种操作分别约束当前求导与跨步变化，不能合并理解。

回看两路输入，目标不是裁剪$B$的像素，而是当前目标网络算出的$\bm{z}^{\mathrm{tar}}_B$。像素可在编码前保留，仍属于第~\ref{sec:data-derived-knowledge-observation}节；表示向量则会随网络变化。它与前章的任务教师输出也不同：分类教师的坐标对应预设类别，BYOL向量坐标没有先命名为花卉物种。名称都叫目标分支或教师，不能代替对输出含义的检查。

预训练结束以后，BYOL通常只保留在线编码器$f_{\bm{\theta}}$的输出供下游使用，舍弃投影器、预测器与目标分支。训练中接近的$\bm{p}_A$和$\bm{z}^{\mathrm{tar}}_B$因而不是直接交给分类器的向量，需要在编码特征上另检验任务效用。

若我们希望匹配的是输出维度上的相对权重，而非归一化连续向量，就需要先把每路输出转成分布。DINO采用这一跨视图目标：学生和教师都有骨干编码器与投影头，结构相同而参数不同，没有BYOL式预测头。投影头为视图$\bm{v}$给出$K$维分数$\bm{u}_{\mathrm{s}}(\bm{v})$与$\bm{u}_{\mathrm{t}}(\bm{v})$，$k,r$枚举输出维度而非物种。为调节教师目标集中在哪些维度，教师先减去按维度维护的中心$\bm{c}$，再与学生分别以正温度$\tau_{\mathrm{t}},\tau_{\mathrm{s}}$归一化：
\begin{align}
  p_{\mathrm{s},k}(\bm{v})
  &=\frac{\exp(u_{\mathrm{s},k}(\bm{v})/\tau_{\mathrm{s}})}
  {\sum_{r=1}^{K}\exp(u_{\mathrm{s},r}(\bm{v})/\tau_{\mathrm{s}})},\\
  p_{\mathrm{t},k}(\bm{v})
  &=\frac{\exp((u_{\mathrm{t},k}(\bm{v})-c_k)/\tau_{\mathrm{t}})}
  {\sum_{r=1}^{K}\exp((u_{\mathrm{t},r}(\bm{v})-c_r)/\tau_{\mathrm{t}})}.
  \label{eq:data-derived-knowledge-dino-distributions}
\end{align}
这些分布给交叉熵提供同一组$K$个比较维度。DINO把视图范围扩展为两个全局视图和若干局部视图：教师只处理两个全局视图，学生处理全局与局部视图，每个教师输出监督除同一视图外的学生输出。沿用正对应的是同图来源，新增要求是局部输入也预测全局分布。教师端停止梯度并以学生参数EMA更新，中心$\bm{c}$另外以教师批量输出的移动平均更新\scite{ka-caron2021dino}

减去中心称为居中（centering），它抑制长期偏高的维度继续支配目标；较低教师温度称为锐化（sharpening），它让维度权重更集中，抵抗所有样本都接近均匀分布的状态。原论文将两者与按EMA更新的动量教师和多裁剪共同使用，并以消融观察作用，不能从其中一个组件推出任意模型都不会坍塌。$K$维仍没有人工类别名称，交叉熵低只说明输出分布匹配，不能认定已发现$K$个真实物种。

DINO的$K$维分布同样只是预训练对象，下游使用骨干特征。本章约定取训练结束时的教师编码器检查点并舍弃投影头；使用学生检查点也是另一种接口选择，需要明确记录，不能把两套编码器的结果混在同一评价中。

这些预测目标省去显式负对照，代价是独立目标分支或更多视图的计算。局部输入缺少全局内容时，能够从输入确定的只是目标的一部分，其余成分需依赖训练分布作估计；两路输出都成为常量时，匹配误差也可能很小。因而需要在更新规则之外检查退化状态，第~\ref{sec:data-derived-knowledge-representation-degeneration}节将把这些状态放回完整目标分析。

\subsubsection{基于对称匹配的目标构造}
\label{sec:data-derived-knowledge-symmetric-matching}

给定同一花朵的两个裁剪，我们也可以不指定哪一路提供固定答案，而同时调整两路使表示接近。\term{对称匹配}（symmetric matching）指同一配对项的梯度同时回到$\bm{z}_i$与$\bm{z}'_i$两路；两路可共享参数，也可使用不同结构，不必先设置独立教师或停止一路梯度。这里对称的是单项求导角色，与BYOL交换两个输入后将非对称损失相加不同。

VICReg把这种要求放在训练嵌入上：编码器$f_{\bm{\theta}}$先取得特征，\term{扩展头}（expander）$h_{\bm{\phi}}$再将特征映射到用于计分的嵌入空间，得到$\bm{z}_i=h_{\bm{\phi}}(f_{\bm{\theta}}(\bm{v}_i))$。本章采用两路共享这两个映射的安排，参数$\bm{\theta},\bm{\phi}$都接收两路梯度。要评价本批$n$对是否接近，将$d$维嵌入按样本行排列为$\bm{Z},\bm{Z}'\in\mathbb{R}^{n\times d}$，$n\geq2$，再平均同一行对应向量的平方距离：
\begin{equation}
  s(\bm{Z},\bm{Z}')
  =\frac{1}{n}\sum_{i=1}^{n}\lVert\bm{z}_i-\bm{z}'_i\rVert_2^2.
  \label{eq:data-derived-knowledge-vicreg-invariance}
\end{equation}
减小式~\eqref{eq:data-derived-knowledge-vicreg-invariance}会让正对应接近，却也允许所有样本输出同一常量而取得零损失。这个反例暴露了缺口：逐对一致没有要求跨样本保留变化。VICReg因此在每一路分别加入方差与协方差约束\scite{ka-bardes2022vicreg}

要检查某一坐标是否对所有样本都给出相同值，需要从逐对距离转到一路批次中该坐标的变化。对第$j$维，先取批内均值$\bar z_j=n^{-1}\sum_i z_{i,j}$，再用分母$n-1$的样本方差计量偏离均值的程度；加上正的数值稳定项$\epsilon$后取平方根，得到用于阈值比较的标准差：
\[
  \begin{aligned}
    \operatorname{Var}(\bm{Z}_{:,j})
    &=\frac{1}{n-1}\sum_{i=1}^{n}(z_{i,j}-\bar z_j)^2,\\
    \sigma_j(\bm{Z})
    &=\sqrt{\operatorname{Var}(\bm{Z}_{:,j})+\epsilon}.
  \end{aligned}
\]
标准差本身只是诊断量，要让训练惩罚变化不足，还需给定正阈值。取$\epsilon>0$且$\gamma>\sqrt{\epsilon}$，使常数批次也低于阈值；对每维只保留$\gamma-\sigma_j$的正部，再对$d$维平均，变化足够的坐标便不受这项惩罚：
\begin{equation}
  v(\bm{Z})
  =\frac{1}{d}\sum_{j=1}^{d}\max(0,\gamma-\sigma_j(\bm{Z})).
  \label{eq:data-derived-knowledge-vicreg-variance}
\end{equation}
式~\eqref{eq:data-derived-knowledge-vicreg-variance}回答哪些维度仍缺少批内变化。$\max(0,\cdot)$选取阈值缺口，而不是寻找一个最坏样本；它与其他项加权优化，属于软惩罚，不能保证每步都达到$\sigma_j\geq\gamma$。

每维都有变化也可能只是在重复同一个变化，例如两个坐标始终都等于$a$。为区分这种冗余，需要计量不同坐标是否一起偏离均值。令$\overline{\bm{z}}=n^{-1}\sum_i\bm{z}_i$，把各样本的中心化向量外积平均，得到$d\times d$样本协方差矩阵：
\[
  \bm{C}(\bm{Z})
  =\frac{1}{n-1}\sum_{i=1}^{n}
  (\bm{z}_i-\overline{\bm{z}})
  (\bm{z}_i-\overline{\bm{z}})^{\top}.
\]
对角元素是已经检查的各维方差，非对角元素才记录两维线性共同变化。为惩罚重复而不压低单维方差，VICReg只累加非对角元素的平方，并按$d$归一化：
\begin{equation}
  c(\bm{Z})
  =\frac{1}{d}\sum_{\substack{j,k=1\\j\neq k}}^{d}C_{j,k}(\bm{Z})^2.
  \label{eq:data-derived-knowledge-vicreg-covariance}
\end{equation}
单独减小协方差也会接受常数输出，所以它要与保留变化和逐对匹配共同使用。将三种要求以正权重合成，训练同时减小正对应距离、标准差缺口和坐标间线性冗余：
\begin{equation}
  \mathcal{L}_{\mathrm{VICReg}}
  =\lambda s(\bm{Z},\bm{Z}')
  +\mu\!\left[v(\bm{Z})+v(\bm{Z}')\right]
  +\nu\!\left[c(\bm{Z})+c(\bm{Z}')\right],
  \label{eq:data-derived-knowledge-vicreg}
\end{equation}
其中$\lambda,\mu,\nu>0$分别控制三种要求的权重。匹配项在对应的样本对上计分，方差和协方差则在每一路整个批次上计算；统计范围从一对扩大到$n$个样本，但没有新增逐对负例标签。

把常数反例代入完整目标：$\bm{Z}=\bm{Z}'$且所有行相同时，匹配与协方差项为零，每路方差惩罚为$\max(0,\gamma-\sqrt{\epsilon})$。上述条件让它为正，教学设定$\gamma=1,\epsilon=10^{-4}$给出$v(\bm{Z})=0.99$；若改为$\gamma=\epsilon=1$，三项却可同时为零。正惩罚也不保证优化逃离该状态，因为精确常数批次的方差梯度仍为零。另一个反例是各行沿$(a,a)$变化：两个坐标都有方差，非零协方差才揭示它们重复同一方向。

这些统计量还依赖批量与嵌入维数。小批量的估计可能波动较大，较大的$d$则增加显式协方差计算和存储。即使变化充分且线性相关较低，变化也可能来自背景或设备，统计惩罚没有赋予物种语义，也没有保证优化收敛。预训练后VICReg舍弃扩展头$h_{\bm{\phi}}$，使用编码器$f_{\bm{\theta}}$的特征，仍需在这个接口上检查下游用途。

现在可以沿图~\ref{fig:data-derived-knowledge-relation-objectives}对照同一对视图的三种计分路径。实线追踪参加损失的前向值，反向虚线追踪当前梯度，点划线另画跨步参数平均。BYOL目标端有前向值而无损失梯度，VICReg同一项却向两路回传；是否再交换输入求和，不会改变这一区别。

\begin{figure*}[htbp]
  \centering
  \begin{tikzpicture}[
  x=1mm,y=1mm,
  font=\figurefont\footnotesize,
  panel/.style={draw=BookRule,fill=NNPanel,line width=.85pt,rounded corners=1.5mm},
  input/.style={draw=NNInput,fill=NNInput!16,line width=1pt,minimum width=7mm,minimum height=7mm,inner sep=1mm,align=center},
  compute/.style={draw=NNHidden,fill=NNHidden!16,line width=1pt,minimum width=10mm,minimum height=8mm,inner sep=1mm,align=center,rounded corners=1mm},
  output/.style={draw=NNOutput,fill=NNOutput!16,line width=1pt,minimum width=8mm,minimum height=8mm,inner sep=1mm,align=center},
  loss/.style={draw=NNGradient,fill=NNGradient!12,line width=1pt,minimum width=10mm,minimum height=8mm,inner sep=1mm,align=center,rounded corners=1mm},
  forward/.style={knowledge arrow,draw=NNWire,line width=1pt},
  gradient/.style={knowledge arrow,draw=NNGradient,line width=1.1pt,dashed,rounded corners=1mm},
  gradient stop/.style={-{Bar[width=2.4mm]},draw=NNGradient,line width=1.1pt,dashed,rounded corners=1mm},
  ema/.style={knowledge arrow,draw=BookTeal,line width=1pt,dash dot},
  note/.style={font=\figurefont\scriptsize,text=BookMuted,align=center}
]
  \draw[panel] (.5,103) rectangle (168.5,153);
  \draw[panel] (.5,48) rectangle (168.5,98);
  \draw[panel] (.5,-10) rectangle (168.5,42);
  \node[font=\figurefont\bfseries\footnotesize,text=BookInk,anchor=west] at (5,148) {(a) 对比判别：SimCLR};
  \node[note,anchor=east] at (163,148) {两路共享编码器与投影器};
  \node[input] (s-v1) at (10,132) {$\bm{v}_1$};
  \node[input] (s-v2) at (10,117) {$\bm{v}_2$};
  \node[compute,minimum width=20mm] (s-f1) at (42,132) {$f_{\bm{\theta}}$};
  \node[compute,minimum width=20mm] (s-f2) at (42,117) {$f_{\bm{\theta}}$};
  \node[compute,minimum width=20mm] (s-g1) at (76,132) {$g_{\bm{\theta}}$};
  \node[compute,minimum width=20mm] (s-g2) at (76,117) {$g_{\bm{\theta}}$};
  \node[output,minimum width=14mm] (s-z1) at (108,132) {$\bm{z}_1$};
  \node[output,minimum width=14mm] (s-z2) at (108,117) {$\bm{z}_2$};
  \node[loss,minimum width=26mm] (s-loss) at (148,124.5) {相似度$\bm{S}$\\与配对损失};
  \draw[forward] (s-v1.east)--(s-f1.west);
  \draw[forward] (s-f1.east)--(s-g1.west);
  \draw[forward] (s-g1.east)--(s-z1.west);
  \draw[forward] (s-v2.east)--(s-f2.west);
  \draw[forward] (s-f2.east)--(s-g2.west);
  \draw[forward] (s-g2.east)--(s-z2.west);
  \draw[NNWire,line width=1pt,rounded corners=1mm] (s-z1.east) -- (126,132) -- (126,124.5);
  \draw[NNWire,line width=1pt,rounded corners=1mm] (s-z2.east) -- (126,117) -- (126,124.5);
  \draw[forward] (126,124.5)--(s-loss.west);
  \draw[gradient] (s-loss.north) -- (148,140) -- (42,140) -- (s-f1.north);
  \draw[gradient] (76,140)--(s-g1.north);
  \draw[gradient] (108,140)--(s-z1.north);
  \draw[gradient] (s-loss.south) -- (148,109) -- (42,109) -- (s-f2.south);
  \draw[gradient] (76,109)--(s-g2.south);
  \draw[gradient] (108,109)--(s-z2.south);

  \node[font=\figurefont\bfseries\footnotesize,text=BookInk,anchor=west] at (5,93) {(b) 跨视图预测：BYOL};
  \node[note,anchor=east] at (163,93) {只画$A\to B$；交换视图另算一项};
  \node[input] (b-va) at (10,78) {$\bm{v}_A$};
  \node[compute,minimum width=20mm] (b-fo) at (42,78) {$f_{\bm{\theta}}$};
  \node[compute,minimum width=20mm] (b-go) at (76,78) {$g_{\bm{\theta}}$};
  \node[compute,minimum width=20mm] (b-qo) at (108,78) {$q_{\bm{\theta}}$};
  \node[input] (b-vb) at (10,59) {$\bm{v}_B$};
  \node[compute,minimum width=20mm] (b-ft) at (42,59) {$f_{\bm{\xi}}$};
  \node[compute,minimum width=20mm] (b-gt) at (76,59) {$g_{\bm{\xi}}$};
  \node[output,minimum width=20mm] (b-sg) at (108,59) {$\operatorname{sg}$};
  \node[loss,minimum width=20mm] (b-loss) at (148,68.5) {$\ell_{A\to B}$};
  \draw[forward] (b-va.east)--(b-fo.west);
  \draw[forward] (b-fo.east)--(b-go.west);
  \draw[forward] (b-go.east)--(b-qo.west);
  \draw[forward,rounded corners=1mm] (b-qo.east) -| (b-loss.north);
  \draw[forward] (b-vb.east)--(b-ft.west);
  \draw[forward] (b-ft.east)--(b-gt.west);
  \draw[forward] (b-gt.east)--(b-sg.west);
  \draw[forward,rounded corners=1mm] (b-sg.east) -| (b-loss.south);
  \draw[gradient] (b-loss.east) -- (164,68.5) -- (164,86) -- (42,86) -- (b-fo.north);
  \draw[gradient] (76,86)--(b-go.north);
  \draw[gradient] (108,86)--(b-qo.north);
  \draw[gradient stop] (164,68.5) -- (164,52) -- (108,52) -- (b-sg.south);
  \draw[ema] (b-fo.south)--(b-ft.north);
  \draw[ema] (b-go.south)--(b-gt.north);
  \node[note,text=BookTeal] at (59,68.5) {EMA};

  \node[font=\figurefont\bfseries\footnotesize,text=BookInk,anchor=west] at (5,37) {(c) 对称匹配：VICReg};
  \node[note,anchor=east] at (163,37) {两路共享编码器与扩展头；无停止梯度};
  \node[input] (v-va) at (10,22) {$\bm{v}_A$};
  \node[compute,minimum width=28mm] (v-fa) at (60,22) {$h_{\bm{\phi}}\circ f_{\bm{\theta}}$};
  \node[output,minimum width=14mm] (v-za) at (100,22) {$\bm{Z}$};
  \node[input] (v-vb) at (10,4) {$\bm{v}_B$};
  \node[compute,minimum width=28mm] (v-fb) at (60,4) {$h_{\bm{\phi}}\circ f_{\bm{\theta}}$};
  \node[output,minimum width=14mm] (v-zb) at (100,4) {$\bm{Z}'$};
  \node[loss,minimum width=34mm] (v-loss) at (148,13) {匹配$s(\bm{Z},\bm{Z}')$\\方差$v(\bm{Z}),v(\bm{Z}')$\\协方差$c(\bm{Z}),c(\bm{Z}')$};
  \draw[forward] (v-va.east)--(v-fa.west);
  \draw[forward] (v-fa.east)--(v-za.west);
  \draw[forward] (v-vb.east)--(v-fb.west);
  \draw[forward] (v-fb.east)--(v-zb.west);
  \draw[NNWire,line width=1pt,rounded corners=1mm] (v-za.east) -- (123,22) -- (123,13);
  \draw[NNWire,line width=1pt,rounded corners=1mm] (v-zb.east) -- (123,4) -- (123,13);
  \draw[forward] (123,13)--(v-loss.west);
  \draw[gradient] (v-loss.north) -- (148,29) -- (60,29) -- (v-fa.north);
  \draw[gradient] (100,29)--(v-za.north);
  \draw[gradient] (v-loss.south) -- (148,-3) -- (60,-3) -- (v-fb.south);
  \draw[gradient] (100,-3)--(v-zb.south);

  \draw[forward] (3,-17) -- (16,-17);
  \node[note,anchor=west] at (18,-17) {前向值};
  \draw[gradient] (69,-17) -- (56,-17);
  \node[note,anchor=west] at (72,-17) {当前步梯度};
  \draw[ema] (120,-17) -- (133,-17);
  \node[note,anchor=west] at (136,-17) {跨步模型参数平均};
\end{tikzpicture}

  \caption{沿三种线型核对关系目标的更新角色。实线传前向值，橙色虚线传当前梯度，青色点划线表示跨步EMA。SimCLR共享编码器与投影器，并使用批内对照；BYOL仅画$A\to B$，sg阻断目标端梯度，EMA只连接编码器和投影器，交换视图需另算一项；VICReg共享$f_{\bm{\theta}}$与$h_{\bm{\phi}}$，匹配项及各路统计项都向相应计算分支回传。}
  \label{fig:data-derived-knowledge-relation-objectives}
\end{figure*}
\FloatBarrier

表~\ref{tab:data-derived-knowledge-relation-objectives}再按比较量、当前梯度和跨步变化核对这些设计。SimCLR使用正对应和负对照，另外三种只用正对应；后者却分别采用独立目标端或两路共同求导。因此无显式负例、有教师、停止梯度与两向求和分别描述不同属性，不能由一项判断其余属性。

\begin{table*}[htbp]
  \centering
  \renewcommand{\arraystretch}{1.15}
  \begin{tabularx}{\linewidth}{@{}>{\RaggedRight\setlength{\parindent}{0pt}\arraybackslash}p{22mm}*{3}{>{\RaggedRight\setlength{\parindent}{0pt}\arraybackslash}X}@{}}
    \toprule
    方法 & 比较对象 & 当前梯度去向 & 跨步变化 \\
    \midrule
    SimCLR & 投影表示的相对相似度 & 两路编码器与投影器 & 共享参数更新后重算表示 \\
    BYOL & 在线预测与目标投影 & 每项仅回传在线分支 & 目标编码器与投影器采用EMA \\
    DINO & 学生与教师的输出分布 & 每项仅回传学生分支 & 教师参数与输出中心分别移动平均 \\
    VICReg & 嵌入匹配与批内统计 & 两路编码器与扩展头 & 共享参数更新后重算表示及统计量 \\
    \bottomrule
  \end{tabularx}
  \caption{固定正对应后，比较量与更新方式仍可不同。BYOL在线端有额外预测器，DINO没有该部件，VICReg则没有独立教师。表中组件承担不同职责，完整退化条件和计算成本需结合正文判断。}
  \label{tab:data-derived-knowledge-relation-objectives}
\end{table*}
\FloatBarrier

要回答这些训练依据是否支持物种识别，可固定任务数据、标签划分和线性分类协议，冻结各方法选定的编码器，用$\bm{h}=f(\bm{x})$拟合分类头。在同一批未参与拟合和调参的任务样本上，按错误类别预测数除以样本数计算误分类率，数值越低表示这一协议下的识别误差越小。此时比较的是编码特征对任务的支持，不是各训练头拟合自身目标的分数。若还更新编码器，比较对象便成为整体微调，应另报协议；投影向量、BYOL预测向量和DINO的$K$维分布也不能混作同一特征接口。

\subsection{目标更新与表示退化的控制}
\label{sec:data-derived-knowledge-target-update}

在BYOL的同一对裁剪上，即使当前目标不接受梯度，下一步重新编码时也可能得到不同向量。要理解这种变化，我们需要依次追踪前向目标、允许的反向路径、接受梯度的参数更新，以及下一步目标的准备。“教师随学生更新”省去了这些时点，无法说明目标何时保持不变、何时又发生变化。

另一个问题是，变化中的目标是否仍保留样本差异。只要两路都给常量，匹配误差也可很小；因此还要检查完整目标施加了哪些区分、方差或去相关要求，哪些退化受到惩罚，哪些仍能停住优化。更新规则回答目标怎样变化，退化检查回答变化中可能丢掉什么，两者不能互相替代。

\subsubsection{目标表示的更新}
\label{sec:data-derived-knowledge-target-representation-update}

先比较下一步怎样取得目标。SimCLR和常见VICReg实现的两路共享参数，一次优化后都用同一套新参数重算；独立目标分支则可全程固定、周期性复制或作移动平均。固定教师意味着全程参数不变，停止梯度只意味着当前求导不沿该路径回传，省略目标计算则连前向值都没有。后两种不能代替第一种，也不能相互代替。

表~\ref{tab:data-derived-knowledge-gradient-and-update}用BYOL式更新检查停止梯度与目标变化如何并存。取教学标量：参与EMA的在线参数$\theta^{(t)}=2.00$、梯度$0.50$、学习率$\eta=0.20$，对应目标参数$\xi^{(t)}=1.60$、系数$m=0.90$。在线参数先变为$2.00-0.20\times0.50=1.90$，目标再按$0.90\times1.60+0.10\times1.90=1.63$改变。目标没有接受这项梯度，仍在跨步平均后改变；在线预测器因没有对应目标参数，不参加这次平均。

\begin{table*}[htbp]
  \centering
  \renewcommand{\arraystretch}{1.15}
  \begin{tabularx}{\linewidth}{@{}>{\RaggedRight\setlength{\parindent}{0pt}\arraybackslash}p{25mm}*{3}{>{\RaggedRight\setlength{\parindent}{0pt}\arraybackslash}X}@{}}
    \toprule
    时点 & 在线分支 & 目标分支 & 本步操作 \\
    \midrule
    1. 前向 & $\bm{\theta}^{(t)}$产生预测 & $\bm{\xi}^{(t)}$产生目标 & 比较预测与停止梯度的目标值 \\
    2. 反向 & 接收梯度；示例为$0.50$ & 不接收损失梯度 & 只沿在线路径累计梯度 \\
    3. 在线更新 & 优化器更新；示例为$1.90$ & 仍为$\bm{\xi}^{(t)}$ & 先完成在线参数更新 \\
    4. 目标更新 & 保持$\bm{\theta}^{(t+1)}$ & EMA更新；示例为$1.63$ & 用式~\eqref{eq:data-derived-knowledge-byol-ema}平均对应参数 \\
    5. 下一步 & 新参数产生新预测 & 新参数重新计算目标 & 进入$t+1$步前向计算 \\
    \bottomrule
  \end{tabularx}
  \caption{按时点区分损失梯度与目标参数变化。固定教师省去第4行，目标参数全程不变；共享参数方案没有独立$\bm{\xi}$，两路都读取同次更新后的参数。VICReg的共同反向路径应按图~\ref{fig:data-derived-knowledge-relation-objectives}理解，不套用独立目标列。}
  \label{tab:data-derived-knowledge-gradient-and-update}
\end{table*}
\FloatBarrier

这个混合速度由$m$控制：较小$m$给新在线参数更大权重，目标跟随较快；$m$接近$1$则保留更多旧值，减缓变化但增加滞后，$m=1$时不再混入新值。实际变化还取决于优化器步长、系数日程和训练长度，单比$m$不能确定结果。独立目标端也要额外保存参数并完成前向计算。

复现这一过程还要保持更新顺序。式~\eqref{eq:data-derived-knowledge-byol-ema}读取的是新在线参数$\bm{\theta}^{(t+1)}$，改读$\bm{\theta}^{(t)}$就增加一步滞后；分布式设备若在EMA之前没有同步在线参数，也可能维护出不同目标网络。记录顺序和同步时点，才能确认下一步计算的是哪一个目标。

在DINO中，教师参数与教师输出中心还各有一次移动平均。前者更新生成特征和分数的网络，后者更新从分数中减去的按维度基准；即使都采用平均，它们保存的状态和作用位置仍不同，需要分别维护。

\subsubsection{表示退化的控制}
\label{sec:data-derived-knowledge-representation-degeneration}

假设两张不同花卉照片及各自裁剪都得到同一向量，配对匹配完全成功，表示却无法区分两张照片。这就是\term{表示坍塌}（representation collapse）的直接形式，所有样本都输出常量。另有\term{维度退化}（dimensional degeneration）：样本虽有变化，多个坐标却重复同一方向，例如$(a,a)$只有一个可独立变化的量，有效维数低于名义维数。只看正对应距离会漏掉这两种情况。

图~\ref{fig:data-derived-knowledge-representation-degeneration}固定四个教学样本与坐标范围，对照常数、单方向变化和二维分散。每个样本的实心点与空心圈都接近，配对误差因而都可很小；比较跨样本方差和坐标协方差，才看得到前两种退化。右侧只排除了这两个示意维度中的简单退化，没有给四个点赋予真实类别。

\begin{figure*}[!t]
  \centering
  \begin{tikzpicture}[
  x=1mm,y=1mm,
  font=\figurefont\scriptsize,
  panel/.style={draw=BookRule,fill=NNPanel,line width=.85pt,rounded corners=1.5mm},
  axis/.style={knowledge arrow,draw=BookMuted,line width=.8pt},
  viewone/.style args={#1}{circle,fill=#1,draw=#1,minimum size=3.2mm,inner sep=0pt},
  viewtwo/.style args={#1}{circle,fill=white,draw=#1,line width=1pt,minimum size=3.8mm,inner sep=0pt},
  samplelabel/.style={font=\figurefont\scriptsize,text=BookInk},
  note/.style={font=\figurefont\scriptsize,text=BookMuted,align=center}
]
  % 所有面板用同一映射：(z_1,z_2) -> (中心横坐标+8z_1,37+8z_2) mm。
  % 坐标范围均为[-2,2]^2；点位直接使用矩阵中的教学数值。
  \draw[panel] (.5,.5) rectangle (51.5,66.5);
  \draw[panel] (58.5,.5) rectangle (110.5,66.5);
  \draw[panel] (117.5,.5) rectangle (168.5,66.5);
  \foreach \cx in {26,84.5,143}{
    \begin{scope}[shift={(\cx,37)},x=8mm,y=8mm]
      \draw[axis] (-2,0) -- (2,0) node[right] {$z_1$};
      \draw[axis] (0,-2) -- (0,2) node[above] {$z_2$};
      \draw[BookMuted,line width=.6pt] (1,-.04) -- (1,.04)
        node[below=.4mm] {$1$};
      \draw[BookMuted,line width=.6pt] (-.04,1) -- (.04,1)
        node[left=.4mm] {$1$};
    \end{scope}
  }
  \node[font=\figurefont\bfseries\footnotesize,text=BookInk] at (26,62) {(a) 常数表示};
  \node[viewtwo=BookAmber,minimum size=7mm] at (26,37) {};
  \node[viewtwo=NNHidden,minimum size=5.8mm] at (26,37) {};
  \node[viewtwo=BookTeal,minimum size=4.6mm] at (26,37) {};
  \node[viewone=BookBlue] at (26,37) {};
  \node[samplelabel,anchor=west] at (30,40) {$A=B=C=D$};
  \node[note] at (26,12) {$\bm{Z}=\bm{Z}'=\begin{bmatrix}0&0\\0&0\\0&0\\0&0\end{bmatrix}$\\配对误差为$0$，各维方差为$0$};

  \node[font=\figurefont\bfseries\footnotesize,text=BookInk] at (84.5,62) {(b) 单方向变化};
  \begin{scope}[shift={(84.5,37)},x=8mm,y=8mm]
    \draw[BookRule,line width=.6pt,dashed] (-1.7,-1.7) -- (1.7,1.7);
    \foreach \x/\y/\sampleid/\samplecolor in {-1.5/-1.5/A/BookBlue,-.5/-.5/B/BookTeal,.5/.5/C/NNHidden,1.5/1.5/D/BookAmber}{
      \node[viewtwo=\samplecolor] at ({\x+.08},{\y-.05}) {};
      \node[viewone=\samplecolor] at (\x,\y) {};
      \node[samplelabel,anchor=south east,xshift=-1mm,yshift=1mm] at (\x,\y) {\sampleid};
    }
  \end{scope}
  \node[note] at (84.5,10.5) {$\bm{Z}\approx\bm{Z}'=\begin{bmatrix}-1.5&-1.5\\-.5&-.5\\.5&.5\\1.5&1.5\end{bmatrix}$\\样本有方差，但两个坐标重复同一变化};

  \node[font=\figurefont\bfseries\footnotesize,text=BookInk] at (143,62) {(c) 示意维度中分散};
  \begin{scope}[shift={(143,37)},x=8mm,y=8mm]
    \foreach \x/\y/\sampleid/\samplecolor in {-1/-1/A/BookBlue,-1/1/B/BookTeal,1/-1/C/NNHidden,1/1/D/BookAmber}{
      \node[viewtwo=\samplecolor] at ({\x+.08},{\y-.05}) {};
      \node[viewone=\samplecolor] at (\x,\y) {};
      \node[samplelabel,anchor=south east,xshift=-1mm,yshift=1mm] at (\x,\y) {\sampleid};
    }
  \end{scope}
  \node[note] at (143,10.5) {$\bm{Z}\approx\bm{Z}'=\begin{bmatrix}-1&-1\\-1&1\\1&-1\\1&1\end{bmatrix}$\\两维有方差且示例协方差为$0$};

  \node[viewone=BookInk] at (5,-5) {};
  \node[note,anchor=west] at (8,-5) {视图A};
  \node[viewtwo=BookInk] at (37,-5) {};
  \node[note,anchor=west] at (40,-5) {对应视图B};
  \node[note,anchor=west,text=FigureLoss] at (76,-5) {教学构造；分散不等于形成有用类别};
\end{tikzpicture}

  \caption{低配对误差没有排除样本或维度退化。相同标记代表同一样本，实心点与空心圈代表其两视图，三个面板使用同一坐标范围。手工矩阵依次给出常数、重复方向与二维分散；分散只描述示意坐标，不代表真实语义类别。}
  \label{fig:data-derived-knowledge-representation-degeneration}
\end{figure*}

对SimCLR，已有计算给出常数输出的损失$\log(2n-1)$：有效对照与可实现的配对优势能使目标更低，精确常数状态却仍可能是驻点。这个目标要求的是当前关系集合中的相对排序，错误负例会排斥语义相近样本，背景差异也可能提供排序优势。因此区分偏好既不是优化逃离常数的保证，也没有直接约束物种语义。

对BYOL，要检查的是预测头、停止梯度、移动平均目标及优化动态共同组成的非对称训练。原论文完整方法没有显式负例，并通过消融研究这些组件\scite{ka-grill2020byol}
停止梯度只阻断当前求导，单靠这一事实不能推出不会坍塌；EMA按旧目标与新在线参数加权，也不是沿目标损失梯度更新。两种操作各自有明确作用，成功边界仍属于完整训练安排。

DINO比较$K$维分布，诊断对象也要相应改变。所有样本都给近均匀权重是一种退化，全部集中到同一维度是另一种；教师锐化抵抗前者，居中抵抗后者，两者在动量教师条件下共同使用\scite{ka-caron2021dino}
它们与连续向量的简单重合不是同一个诊断口径，需要查看输出分布，不能仅凭一张散点图排除。

VICReg的附加计分可以用图中矩阵直接核对。在$n\geq2$、$\gamma>\sqrt{\epsilon}$的条件下，方差项惩罚(a)的常数状态，协方差项惩罚(b)的重复方向。用分母$n-1$计算，(b)的协方差矩阵各元素为$5/3$，非对角平方惩罚为$c(\bm{Z})=25/9$；(c)的矩阵为$(4/3)\bm{I}_2$，$\bm{I}_2$是二维单位矩阵，故$c(\bm{Z})=0$。这验证计分能区分坐标冗余，没有证明任意初值下都能避开退化。

排除这些退化后，还要回到花卉任务：较高方差可能只来自照明，低协方差也可能只是背景特征去相关，分散并没有证明花瓣信息被保留。至少应分别观察匹配误差、每维方差、协方差或有效秩，以及独立下游任务的指标；有效秩概括表示实际利用的变化方向有多少。前几项检查差异是否存在，最后一项检查差异是否有用，训练总损失不能替代它们。

\section{基于数据分组的知识获取}
\label{sec:data-derived-knowledge-grouping}

若要让一批花卉照片共享群组目标，仅记录若干样本对还不足以决定整个分组。我们可以先根据表示的相似性组织群组，再让模型预测所属组号或分配权重。\term{分组目标}（grouping target）就是由当前表示、比较规则与分组配置生成的临时答案。范围从样本对扩大到一批样本，目标身份也变为算法产生的组号；它没有预设物种含义，下一轮的组1也未必对应这一轮的组1。

聚类的结构假设和求解方法见\BookChapterRef{chap:clustering}{卷4《模型：经典模型》第4章“聚类”}。在这里，已有观测先经表示映射，再由分组规则产生目标，模型读取原输入预测这个目标。需要就地说明的是分组用哪种表示、产生硬编号还是软分配、何时重算，以及重算后怎样同步预测接口，而不把簇存在误当成已有类别答案。

图~\ref{fig:data-derived-knowledge-grouping-targets}固定同一批表示，对照整体聚类、代表向量分配与图连接怎样产出目标。三路可共用编码器，也可组合：原型分配决定目标权重，跨视图交换则沿用第~\ref{sec:data-derived-knowledge-relations}节的正对应。这里按主要目标构造机制归类，不把使用关系的方法排除在分组之外。

\begin{figure*}[!t]
  \centering
  \begin{tikzpicture}[
  x=1mm,y=1mm,
  font=\figurefont\scriptsize,
  panel/.style={draw=BookRule,fill=NNPanel,line width=.85pt,rounded corners=1.5mm},
  source/.style={draw=FigureInput,fill=FigureInput!12,line width=1.1pt,minimum height=9mm,minimum width=28mm,align=center,inner sep=1.5mm},
  target/.style={draw=FigureOutput,fill=FigureOutput!10,line width=1.1pt,minimum height=8mm,minimum width=23mm,align=center,inner sep=1.5mm},
  model/.style={draw=FigureModel,fill=FigureModel!12,line width=1.1pt,minimum height=8mm,minimum width=23mm,align=center,inner sep=1.5mm},
  flow/.style={knowledge arrow,draw=NNWire,line width=1.1pt},
  refresh/.style={knowledge arrow,draw=BookTeal,line width=.9pt,dashed,rounded corners=1.5mm,preaction={draw=white,line width=2.5pt,solid}},
  edge/.style={draw=NNLine,line width=.8pt},
  warningedge/.style={draw=FigureLoss,line width=1.2pt,dashed},
  point/.style={circle,draw=BookBlue,fill=BookBlue!18,minimum size=3.4mm,inner sep=0pt},
  note/.style={font=\figurefont\scriptsize,text=BookMuted,align=center}
]
  \draw[panel] (.5,-34) rectangle (52.5,34.5);
  \draw[panel] (58.5,-34) rectangle (110.5,34.5);
  \draw[panel] (116.5,-34) rectangle (168.5,34.5);
  \node[source] (batch) at (84.5,57) {同一批未标注样本$\bm{X}$};
  \node[model] (repr) at (84.5,44) {共享编码器$f_{\bm{\theta}}$\\$\bm{Z}=f_{\bm{\theta}}(\bm{X})$};
  \draw[flow] (batch)--(repr);
  \draw[NNWire,line width=1.1pt] (repr.south) -- (84.5,38);
  \draw[NNWire,line width=1.1pt] (5,38) -- (164,38);

  \node[font=\figurefont\bfseries\footnotesize,text=BookInk] at (26.5,30) {(a) 聚类编号};
  \draw[flow] (5,38) |- (26.5,24);
  \foreach \x/\y in {10/16,14/20,18/14,34/17,39/21,43/14}{\node[point] at (\x,\y) {};}
  \draw[BookBlue,line width=.8pt,dashed] (14,17) ellipse (9mm and 7mm);
  \draw[BookTeal,line width=.8pt,dashed] (39,17) ellipse (9mm and 7mm);
  \node[text=BookBlue] at (14,8) {组1};
  \node[text=BookTeal] at (39,8) {组2};
  \node[target] (cluster-target) at (26.5,-5) {临时目标$c_i$};
  \node[model] (cluster-predict) at (26.5,-24) {预测组号并更新};
  \draw[flow] (26.5,7)--(cluster-target);
  \draw[flow] (cluster-target)--(cluster-predict);
  \draw[refresh] (cluster-predict.east) -- (54,-24) |- (repr.west);
  \node[note,text=BookTeal,rotate=90] at (47,-9) {可选重编码};

  \node[font=\figurefont\bfseries\footnotesize,text=BookInk] at (84.5,30) {(b) 原型软分配};
  \node[note,text=BookInk] (assignments) at (84.5,15) {$\begin{array}{c|cc}
    &\text{样本1}&\text{样本2}\\
    \bm{c}_1&0.690&0.310\\
    \bm{c}_2&0.310&0.690
  \end{array}$};
  \draw[flow] (61,38) |- (assignments.west);
  \node[target] (proto-target) at (84.5,-5) {分配向量$\bm{q}_{s,b}$};
  \node[model] (proto-predict) at (84.5,-24) {视图$t$预测并更新};
  \draw[flow] (assignments.south)--(proto-target.north);
  \draw[flow] (proto-target)--(proto-predict);
  \draw[refresh] (proto-predict.east) -- (112,-24) |- (repr.east);

  \node[font=\figurefont\bfseries\footnotesize,text=BookInk] at (142.5,30) {(c) 邻域／图分组};
  \draw[flow] (119,38) |- (142.5,26.5);
  \node[point,minimum size=4.3mm] (g1) at (125,17) {1};
  \node[point,minimum size=4.3mm] (g2) at (132,23) {2};
  \node[point,minimum size=4.3mm] (g3) at (136,12) {3};
  \node[point,minimum size=4.3mm] (g4) at (150,12) {4};
  \node[point,minimum size=4.3mm] (g5) at (156,23) {5};
  \node[point,minimum size=4.3mm] (g6) at (163,17) {6};
  \draw[edge] (g1)--(g2);
  \draw[edge] (g2)--(g3);
  \draw[edge] (g3)--(g1);
  \draw[edge] (g4)--(g5);
  \draw[edge] (g5)--(g6);
  \draw[edge] (g6)--(g4);
  \draw[warningedge] (g3.east)-- node[above=1mm,text=FigureLoss] {桥接} (g4.west);
  \node[target,text width=38mm,inner sep=1mm] (graph-target) at (142.5,-5) {前：$(1,1,1,2,2,2)$\\后：$(1,1,1,1,1,1)$};
  \node[model] (graph-predict) at (142.5,-24) {预测目标并更新};
  \draw[flow] (142.5,7)--(graph-target);
  \draw[flow] (graph-target)--(graph-predict);
  \draw[refresh] (graph-predict.east) -- (167,-24) |- (repr.north east);

  \node[note,anchor=north west,text width=166mm,align=left] at (1,-38) {原型面板只展开两个样本，列内权重和为1；图面板展开六个节点。数值均为教学构造，分组和原型索引不具有预设物种含义。};
\end{tikzpicture}

  \caption{由同批表示生成三种分组目标。聚类给临时组号，原型面板的两列为第~\ref{sec:data-derived-knowledge-prototype-targets}节算例的归一化分配，图面板六个节点用于第~\ref{sec:data-derived-knowledge-graph-targets}节的连通计算；两列与六节点不是同一展开规模。赭色虚线表示拟加入的错误边$(3,4)$，青色回线表示更新后重算目标，不表示向分组求解器回传梯度。}
  \label{fig:data-derived-knowledge-grouping-targets}
\end{figure*}

表~\ref{tab:data-derived-knowledge-grouping}按输出、统计范围和刷新状态读这三路。全量聚类、批内分配与邻域建图所需的信息范围不同；重算分组后，与之相连的编号、原型或图关系及预测接口也要同步检查，只换目标文件可能留下旧含义的输出头。

\begin{table*}[htbp]
  \centering
  \renewcommand{\arraystretch}{1.15}
  \begin{tabularx}{\linewidth}{@{}>{\RaggedRight\setlength{\parindent}{0pt}\arraybackslash}p{25mm}*{3}{>{\RaggedRight\setlength{\parindent}{0pt}\arraybackslash}X}@{}}
    \toprule
    分组方式 & 输出目标 & 统计范围 & 刷新时同步对象 \\
    \midrule
    聚类分组 & 硬簇编号$c_i$ & 常需全数据或较大样本 & 簇中心、样本目标、预测头与编号对应 \\
    原型分配 & 硬编号或软分配$\bm{q}_i$ & 小批量，可辅以表示队列 & 原型、分配结果与队列 \\
    邻域／图分组 & 组号或组内关系 & 局部近邻或全局图 & 邻接表、分组结果与训练目标 \\
    \bottomrule
  \end{tabularx}
  \caption{比较分组输出及其刷新所依赖的状态。三种机制可以组合，硬编号和软权重都由算法构造，没有预设类别含义；统计范围和同步对象不同，不能用一套刷新规则代替全部机制。}
  \label{tab:data-derived-knowledge-grouping}
\end{table*}

% 首段的DeepCluster边注从段末开始，避免贴近纸张底边。
\Needspace{65mm}
\subsection{基于聚类的分组目标构造}
\label{sec:data-derived-knowledge-cluster-targets}

对未标注物种的照片，先按当前特征分组，就能让每张照片暂时获得一个组号。具体地，提取$\bm{z}_i=f_{\bm{\theta}}(\bm{x}_i)$，按所选聚类规则得到$c_i\in\{1,\ldots,K\}$，再从原输入训练模型预测$c_i$。DeepCluster以$k$均值分组网络特征，并用簇分配监督网络更新，给出了这种交替机制的代表实现\scite{knowledge-acquisition-deepcluster2018}

组号可以接到分类式损失，但分类形式没有使它变成物种标签：编号由当前表示和聚类配置决定。以下把这一类交替训练抽象成提取、分组与预测三步；跨轮编号对齐、刷新策略和失效检查是使用这种机制还需处理的问题，不全部归因于该代表论文。

训练开始前还需选择答案是否刷新。固定分组只计算一次$c_i$，之后始终预测这些编号；周期性分组则要记录每轮使用的参数、编号和训练集，才能分开“目标生成”与“拟合目标”。对固定的$n$个样本，交替过程写为
\[
  \begin{aligned}
    \{\bm{z}_i^{(r)}=f_{\bm{\theta}^{(r)}}(\bm{x}_i)\}_{i=1}^{n}
    &\ \longrightarrow\
    \bm{c}^{(r)}
    =\operatorname{Cluster}(\{\bm{z}_i^{(r)}\}_{i=1}^{n}),\\
    \mathcal{D}^{(r)}=\{(\bm{x}_i,c_i^{(r)})\}_{i=1}^{n}
    &\ \longrightarrow\
    \bm{\theta}^{(r+1)}=\operatorname{Train}(\mathcal{D}^{(r)}),
  \end{aligned}
\]
其中$\bm{c}^{(r)}=(c_i^{(r)})_{i=1}^{n}$，$r$编号的是分组轮次，不是单个梯度步；$\operatorname{Cluster}$按当前特征生成编号，$\operatorname{Train}$表示据本轮输入目标对拟合网络的过程。前一轮训练改变表示，下一轮分组就可能改写答案。频繁刷新缩短目标相对表示的滞后，却增加全量编码、聚类和目标存储。

刷新时不能把数字本身当作群组身份。即使几何分组不变，原组1也可能被算法改名为组3；跨轮保留预测头，就需按样本重叠或中心匹配对齐编号，或重新初始化相应输出。分裂与合并时可能没有一一对应，此时更不能只因数字相同便继承旧参数。

照片可能按背景颜色聚组，也可能按花瓣形态聚组，两种划分都可能符合各自采用的聚类计分，却给物种识别提供不同信息。簇数、特征尺度、异常点、组间不平衡与初始表示偏向都会改变答案；模型再去预测这些答案，还可能强化不合任务需要的分组。

因此，低聚类目标或准确预测簇号，只说明形成或复现了当前分组。把某簇解释为任务类别，需要额外标签、领域解释、独立数据稳定性或下游用途提供相应证据；取得这些证据之前，编号仍是临时目标。

\subsection{基于原型的分配目标构造}
\label{sec:data-derived-knowledge-prototype-targets}

给一张照片分组时，也可以先用少量代表向量与它的表示比较，而不直接维护所有样本的群组边界。\term{原型}（prototype）就是这些可更新的代表$\bm{c}_1,\ldots,\bm{c}_K$。对$\bm{z}_i$，最近原型可给出硬编号，多个原型的支持权重则形成软分配$\bm{q}_i\in\mathbb{R}^{K}$，各权重非负且和为$1$。这些权重描述算法对代表的分配，不会自动成为物种类别的后验概率。

SwAV进一步要求同一照片的不同裁剪能预测彼此的原型分配，沿用正对应而把答案改为分组权重。先固定单位长度的表示和原型，以内积计量支持，再在$K$个原型上指数化、归一化，得到可放进交叉熵的预测分布。对同图视图$t,s$，归一化表示记为$\bm{z}_t,\bm{z}_s$，待构造的分配为$\bm{q}_t,\bm{q}_s$；取温度$\tau>0$，视图表示$\bm{z}$预测原型$k$的概率为
\[
  p_k(\bm{z})
  =\frac{\exp(\bm{z}^{\top}\bm{c}_k/\tau)}
  {\sum_{r=1}^{K}\exp(\bm{z}^{\top}\bm{c}_r/\tau)}.
\]
这里$p_k$是模型的预测，还没有规定目标分配怎样产生。先把另一视图的归一化分配视为当步答案，按$K$个原型的目标权重计算负对数计分；在目标固定时减小交叉熵，就是让预测分布更接近该分配，而非要求每个原型概率同时增大。交换两侧相加，得到
\begin{equation}
  \mathcal{L}_{\mathrm{swap}}
  =-\sum_{k=1}^{K}q_{s,k}\log p_k(\bm{z}_t)
   -\sum_{k=1}^{K}q_{t,k}\log p_k(\bm{z}_s).
  \label{eq:data-derived-knowledge-swav-swap}
\end{equation}
式~\eqref{eq:data-derived-knowledge-swav-swap}中$\bm{q}_s$从$s$侧生成，却核对$t$侧预测，另一项反向使用同一关系；两项采用求和而非方向平均\scite{ka-caron2020swav}

怎样构造这个答案，是原型分配的关键。若每个样本独立选择最近代表，全部样本可能集中到一个原型；SwAV在同一分配批次内要求原型分担总权重，同时保留相似度支持。于是需要一个记录“各样本向各原型分配多少”的非负矩阵，并在满足行列总量的候选矩阵中联合选择，而不是逐样本取最大分数。对正整数$B,K$，令$\bm{Z}=[\bm{z}_1,\ldots,\bm{z}_B]\in\mathbb{R}^{d\times B}$、$\bm{C}=[\bm{c}_1,\ldots,\bm{c}_K]\in\mathbb{R}^{d\times K}$；这里按列排样本，与VICReg按行排批次不同，$\bm{C}$也改指原型矩阵而非协方差。内积加权总分越大越匹配；熵正则（entropy regularization）则额外奖励不过分集中的软权重，使相似度支持与分配集中程度共同参与选择。SwAV在允许集合$\mathcal{Q}$内最大化二者之和：
\begin{align}
  \max_{\bm{Q}\in\mathcal{Q}}\quad&
  \operatorname{Tr}(\bm{Q}^{\top}\bm{C}^{\top}\bm{Z})
  +\varepsilon H(\bm{Q}),\\
  \mathcal{Q}
  ={}&
  \left\{
    \bm{Q}\in\mathbb{R}_{+}^{K\times B}
    \ \middle|\
    \begin{aligned}
      \bm{Q}\bm{1}_B&=\frac{1}{K}\bm{1}_K,\\
      \bm{Q}^{\top}\bm{1}_K&=\frac{1}{B}\bm{1}_B
    \end{aligned}
  \right\}.
  \label{eq:data-derived-knowledge-swav-assignment}
\end{align}
这里$\operatorname{Tr}(\bm{Q}^{\top}\bm{C}^{\top}\bm{Z})$等于$\sum_{k,b}Q_{k,b}\bm{c}_k^{\top}\bm{z}_b$，直接汇总分配与相似度的对应；$H(\bm{Q})=-\sum_{k=1}^{K}\sum_{b=1}^{B}Q_{k,b}\log Q_{k,b}$为熵，约定$0\log0=0$，$\varepsilon>0$控制其权重。$\bm{1}_B,\bm{1}_K$为相应维数的全一向量。集合要求每原型行和为$1/K$、每样本列和为$1/B$，总量为$1$，这是人为施加的均衡条件，不是观测出的类别频率。数学约束为等式，数值求解近似满足；列和还不是概率分布所需的$1$，所以进入交换损失前必须归一化：
\begin{equation}
  \bm{q}_b
  =
  \frac{\bm{Q}_{:,b}}
       {\bm{1}_K^{\top}\bm{Q}_{:,b}}
  =
  B\bm{Q}_{:,b},
  \qquad
  \bm{1}_K^{\top}\bm{q}_b=1.
  \label{eq:data-derived-knowledge-swav-code-normalization}
\end{equation}
式~\eqref{eq:data-derived-knowledge-swav-code-normalization}把第$b$列变成和为$1$的\term{分配向量}（assignment vector），即该样本对$K$个原型的目标权重。交换损失中的$\bm{q}_t,\bm{q}_s$始终指这种向量，原论文称为code，并非程序代码。若直接用和为$1/B$的原始列，交叉熵会额外缩小$1/B$，计分就不再是上述归一化目标。

联合分配还需实际求得可用矩阵。对分数$\bm{F}=\bm{C}^{\top}\bm{Z}$，熵正则化使分配可从正的指数权重开始，再通过缩放满足边际总量。先逐元素指数化并把总和归一化为$1$：
\[
  Q_{k,b}^{(0)}
  =\frac{\exp(F_{k,b}/\varepsilon)}
  {\sum_{r=1}^{K}\sum_{a=1}^{B}\exp(F_{r,a}/\varepsilon)}.
\]
这个初值总和正确，行列边际却未必正确。Sinkhorn--Knopp迭代通过交替缩放修正它们：先把每行总量调到$1/K$，再把每列总量调到$1/B$。用$t$编号这里的缩放轮次，在正权重条件下分母非零，每轮写为
\begin{equation}
  \begin{aligned}
    R_{k,b}^{(t)}
    &=\frac{Q_{k,b}^{(t-1)}}{K\sum_{a=1}^{B}Q_{k,a}^{(t-1)}},\\
    Q_{k,b}^{(t)}
    &=\frac{R_{k,b}^{(t)}}{B\sum_{r=1}^{K}R_{r,b}^{(t)}}.
  \end{aligned}
  \label{eq:data-derived-knowledge-sinkhorn-scaling}
\end{equation}
列缩放可能破坏刚得到的行和，因此需要重复，并用两类边际偏差检查停止。下面教学算例约定：列缩放后，全部行偏离$1/K$和全部列偏离$1/B$的最大绝对值小于$10^{-10}$时停止。这一最大值覆盖每个边际，避免平均误差隐藏某行不达标；列和在数值精度内为$1/B$后，可用$\bm{q}_b=B\bm{Q}_{:,b}$。工程实现若固定少量轮次，还需报告残余行边际误差。

用两个样本和两个原型，我们可以从分数一直复算到交换损失。以下全是教学构造，不是训练实验：$K=B=2$，单位原型$\bm{c}_1=(1,0)^{\top}$、$\bm{c}_2=(0,1)^{\top}$，视图$s$的单位表示为$\bm{z}_{s,1}=(1,0)^{\top}$、$\bm{z}_{s,2}=(0.8,0.6)^{\top}$，$\varepsilon=\tau=0.5$。据内积和指数初始化，得到
\[
  \bm{F}_s=\begin{bmatrix}1&0.8\\0&0.6\end{bmatrix},
  \qquad
  \bm{Q}^{(0)}\approx
  \begin{bmatrix}0.443462&0.297262\\0.060016&0.199260\end{bmatrix}.
\]
将初值各行先调到$0.5$，再将各列调到$0.5$，首轮结果分别为
\[
  \begin{aligned}
    \bm{R}^{(1)}&\approx
    \begin{bmatrix}0.299344&0.200656\\0.115738&0.384262\end{bmatrix},\\
    \bm{Q}^{(1)}&\approx
    \begin{bmatrix}0.360584&0.171525\\0.139416&0.328475\end{bmatrix}.
  \end{aligned}
\]
行缩放后的$\bm{R}^{(1)}$列和约为$0.415081,0.584919$，尚不均衡；列缩放后的$\bm{Q}^{(1)}$行和又变成约$0.532109,0.467891$。继续用未舍入值交替缩放，12轮后的最大边际误差约为$1.82\times10^{-11}$，达到本例停止条件，归一化列后得到
\[
  \bm{Q}_s\approx
  \begin{bmatrix}0.344987&0.155013\\0.155013&0.344987\end{bmatrix},
  \qquad
  \bm{q}_{s,1}\approx\begin{bmatrix}0.689974\\0.310026\end{bmatrix},
  \quad
  \bm{q}_{s,2}\approx\begin{bmatrix}0.310026\\0.689974\end{bmatrix}.
\]
两个样本独立按最大内积分配，都会选原型1；联合均衡分配却使样本2更多支持原型2。这不是相似度被算错，而是允许集合加入了批次均衡要求。一个样本的目标因而依赖同批其他样本，不能仅从自身最近原型还原。

保持样本1、2的正对应，再构造$t$侧表示$\bm{z}_{t,1}=(0.8,0.6)^{\top}$、$\bm{z}_{t,2}=(1,0)^{\top}$。数值位置虽然交换，样本身份没有交换，不能据此重配两侧。对$t$侧独立运行同一分配，得到$\bm{q}_{t,1}=\bm{q}_{s,2}$、$\bm{q}_{t,2}=\bm{q}_{s,1}$；模型对样本1的两侧预测概率为
\[
  \bm{p}(\bm{z}_{t,1})\approx(0.598688,0.401312)^{\top},
  \qquad
  \bm{p}(\bm{z}_{s,1})\approx(0.880797,0.119203)^{\top}.
\]
用另一侧分配核对这些概率，并采用自然对数，样本1的两向计分及其和为
\[
  \begin{aligned}
    -\bm{q}_{s,1}^{\top}\log\bm{p}(\bm{z}_{t,1})&\approx0.637025,\\
    -\bm{q}_{t,1}^{\top}\log\bm{p}(\bm{z}_{s,1})&\approx1.506877,\\
    \mathcal{L}_{\mathrm{swap},1}&\approx2.143902.
  \end{aligned}
\]
样本2的两项在本构造中恰好交换，故对样本平均、对方向求和仍为$2.143902$；再平均两个方向则变成$1.071951$。两者只因聚合口径不同便相差一倍，比较损失前必须固定方向和样本的求和或平均范围。

目标求解和网络更新也有不同路径。当步把Sinkhorn得到的$\bm{q}_t,\bm{q}_s$停止求导，编码器与原型只接收式~\eqref{eq:data-derived-knowledge-swav-swap}中预测$p_k(\bm{z})$一侧的梯度。原方法在优化更新后把非零原型列投影回单位长度，$\bm{c}_k\leftarrow\bm{c}_k/\lVert\bm{c}_k\rVert_2$，下一批再计算分数；这是更新后的约束操作，并非对分配目标反向传播。SwAV没有BYOL式目标网络，让梯度穿过分配求解器或省略原型归一化都会改变原方法的计算路径。

我们还要区分“存在允许分配”与“分配有信息”。任意正整数$K,B$下，$Q_{k,b}=1/(KB)$都满足全部边际，即使表示完全相同，软约束集合也不为空。$K>1$时均衡排除所有权重集中到一个原型，却接受每样本均匀分配到全部原型。真正困难是有信息的近硬分配：多样性不足时，均衡可能迫使相似样本分摊，或与相似度支持的分组冲突。原论文用前批表示队列扩大统计范围，补充可能缺少的多样性，不是修补可行性；队列同时带来旧表示问题，过大的$\varepsilon$也会进一步偏好均匀权重。

目标还随原型数、相似度温度、熵权重、均衡要求和批量组成改变。比较最近原型的硬编号与软分配，应固定原型及表示，才知道差别来自分配规则：硬化删除多原型支持，软权重增加存储和交叉熵计算。本文SwAV采用上面的严格边际要求，其他实现若放宽均衡应另行说明；任何编号都没有预设物种含义。

至此，SwAV的两种依据可以分别追溯：同图正对应来自数据变换，分配答案来自带均衡条件的原型求解。非对比匹配描述它怎样使用关系，原型分组描述它怎样产生目标；只用其中一个名称，都不能完整说明这条训练路径。

\subsection{基于邻域与图结构的分组目标构造}
\label{sec:data-derived-knowledge-graph-targets}

对于一批花卉照片，我们也可以先记录每张照片与哪些照片相近，再从连接中形成群组，而不用少量原型代表全部样本。按$k$近邻、半径阈值或已有关系建立图$\mathcal{G}=(V,E)$，顶点对应样本，边表示规则认定的邻近；这里$V$改指样本顶点，不再是观测位置集合。直接使用邻居集合得到局部目标，采用连通分量、谱划分或社区发现则把局部联系进一步汇成分组编号。

从局部边变成全组目标时，需要检查分组采用的传递规则。$A$接近$B$、$B$接近$C$没有证明$A$与$C$共享语义，连通分量却会沿路径将它们归为同组。图~\ref{fig:data-derived-knowledge-grouping-targets}(c)的赭色桥接边因此可能合并两个密集群组，单边错误的影响不再局限于这一对端点。

固定规则可以复算这种扩大。用图中六个教学节点$V=\{1,\ldots,6\}$，对不同节点在对称相似度不低于$0.8$时连无向边，再按连通分量分组，以每分量最小节点的顺序命名。设$\{1,2,3\}$与$\{4,5,6\}$各自内部两两分数为$0.9$，跨组为$0.2$，规则产生
\[
  \begin{aligned}
    E_{\mathrm{before}}
    &=\{(1,2),(1,3),(2,3),(4,5),(4,6),(5,6)\},\\
    \bm{c}_{\mathrm{before}}&=(1,1,1,2,2,2).
  \end{aligned}
\]
只把节点3与4的分数错误改为$0.85$，其余不变，新边$(3,4)$就连通全部节点，目标变为$\bm{c}_{\mathrm{after}}=(1,1,1,1,1,1)$。若再把同组不同节点的无序对全作为正对应，原来两组共6对，现在共15对，新增9对跨原组关系。记录$(\bm{x}_4,c_4)$从目标2改为1，节点1与6也成为正对应，直接分数却仍为$0.2$。教学分数说明连通规则如何放大关系范围，没有证明六个对象真实同类。

换用其他规则时，需重新明确节点、度量、邻居数或阈值、方向、边权及最终划分，不能继承算例中全部连通结论。固定$k$可能让稀疏区域连接远样本，固定半径又可能让密集区域产生过多边。互近邻、局部尺度或密度规则能调整这些行为，选择仍依赖数据分布，没有统一最优配置。

刷新图分组要重新编码、检索近邻、建边和划分。对$n$个样本直接全对比较，精确检索需$O(n^2)$次距离计算；显式保存全部两两距离才需要$O(n^2)$存储，仅保存每点$k$个近邻的最终邻接表为$O(nk)$，检索索引另计。近似检索可降低计算，却可能漏边或加错边。重算后编号也可能置换、分裂或合并，训练目标需同步刷新，无法确定的关系仍应保留。

把群组再转成正对应，可以接入第~\ref{sec:data-derived-knowledge-relations}节的关系目标，但新关系继承的是分组假设，而非额外确认的语义。这里先从数据关联发现群组，没有类别锚点；前章利用已有任务标签传播答案则有另一种来源依据。两者都可用图，只有额外证据才能把本节群组解释为任务类别。

图刷新还形成一个反馈回路：表示决定近邻，近邻目标再拉近这些表示，稳定的错误边就可能持续强化。保留边的来源、与阈值的距离裕量及跨轮稳定性，能帮助追查这类错误；独立用途检验再回答连接是否有用，图损失下降不能独自完成这两项判断。

\section{本章小结}
\label{sec:data-derived-knowledge-summary}

回到没有物种标签的花卉照片，我们已经能构造三类训练依据。保留花瓣像素或后续影像，就有模型计算之外的内容答案；记录同对象视图、自然配对或负对照，就有表示该共享或区分什么的关系；用表示形成簇、原型分配或图群组，就有算法生成的临时目标。它们都能支持预测，却分别依赖保留观测、成立的关系和分组规则，不能从损失形式推断来源。

要使这些依据可检查，观测构造必须隔离答案并固定损失位置；关系目标必须说明比较对象、候选集合及梯度路径。对比判别在候选中奖励配对优势，跨视图预测以向量或输出分布为当步答案，对称匹配则同时调整两路并检查跨样本变化。停止梯度只约束当步求导，目标可在以后因共享参数、复制或EMA改变；任务教师的预设类别与表示目标的未命名坐标也不能混用。

分组还要求同步维护目标身份。聚类编号可跨轮置换，原型分配受批次均衡与熵权重影响，图连通又会扩大局部错误边的后果；这些组号都尚未等同于物种。SwAV把原型答案与跨视图关系联合使用，说明来源和利用方式可以交叉，联合目标仍需分别说明依据、更新顺序与兼容条件。

因此本章回答的是怎样从已有数据建立训练依据，没有保证这些依据自动成为可靠的物种知识。遮蔽比例、变换、采样、目标更新、统计要求、原型数和建图尺度都加入假设并消耗计算；精确常数驻点、均匀分配和背景捷径仍是需要检查的边界。评价应分开确认构造条件、表示差异与下游效用。第~\ref{chap:knowledge-fusion}章继续组织多种来源的共同使用，第~\ref{chap:knowledge-quality-assessment-correction}章再把这些不同证据接成可复核的质量协议。
